\chapter{Complexity — Step Block}
%%% generated by the blueprint pipeline from TCSlib.Complexity.TuringMachine.Build.StepBlock
\section{Overview}

Declarations of the Lean module \texttt{TCSlib.Complexity.TuringMachine.Build.StepBlock}.

\section{Declarations}

\begin{definition}[Work tape read configurations]
\lean{Turing.workReadsK}
\label{Turing.workReadsK}
\leanok
For each natural number $k$, define a list of all assignments of a possibly blank
Boolean symbol to each of $k$ work tape heads. For $k=0$ this list consists of the
unique empty assignment; for $k+1$, it is formed by choosing the first head’s value from
blank, false, or true and adjoining each assignment for the remaining $k$ heads.
\end{definition}

\begin{definition}[Transition table bit list]
\lean{Turing.tableBitsK}
\label{Turing.tableBitsK}
\leanok
\uses{Turing.Action, Turing.actionBitsK, Turing.workReadsK}
Let $k,n\in\mathbb{N}$ and let a transition rule assign an action to each of $n+1$
states, each possibly blank Boolean input symbol, and each configuration of the $k$ work
tape reads. The table bit list concatenates, in order, the bit lists encoding the
actions for every such combination, with states and input symbols enumerated in the
given order and work tape reads enumerated by the recursive list described above.
\end{definition}

\begin{definition}[Output summary]
\lean{Turing.StepBlock.outputSummary}
\label{Turing.StepBlock.outputSummary}
\leanok
A list of Boolean symbols is summarized by no output if it is empty, by true if it
consists of the single symbol true, and by false in every other case.
\end{definition}

\begin{definition}[Summary update from an emitted symbol]
\lean{Turing.StepBlock.summaryEmit}
\label{Turing.StepBlock.summaryEmit}
\leanok
Given an existing optional Boolean summary and an optional emitted Boolean symbol,
retain the existing summary if no symbol is emitted. If a symbol is emitted, the result
is true exactly when the prior summary is absent and the emitted symbol is true; in all
other emission cases the result is false.
\end{definition}

\begin{lemma}[Summary of an Appended Output Symbol]
\lean{Turing.StepBlock.outputSummary_append}
\label{Turing.StepBlock.outputSummary_append}
\leanok
\uses{Turing.StepBlock.outputSummary, Turing.StepBlock.summaryEmit}
\difficulty{4}
For every finite binary word $w$ and optional bit $b$, the summary of the word formed by
appending $b$ (when present) equals the result of updating the summary of $w$ by
emitting $b$.
\end{lemma}

\begin{lemma}[Characterization of a True Output Summary]
\lean{Turing.StepBlock.outputSummary_true}
\label{Turing.StepBlock.outputSummary_true}
\leanok
\uses{Turing.StepBlock.outputSummary}
\difficulty{4}
For every finite binary word $w$, its output summary is the present bit $1$ if and only
if $w$ consists of exactly the single bit $1$.
\end{lemma}

\begin{lemma}[Emitting from a false summary]
\lean{Turing.StepBlock.summaryEmit_other}
\label{Turing.StepBlock.summaryEmit_other}
\leanok
\uses{Turing.StepBlock.summaryEmit}
\difficulty{3}
For every $b\in\mathrm{Option}(\mathrm{Bool})$, combining the summary value false with
$b$ returns the summary value false.
\end{lemma}

\begin{lemma}[Emitting from an empty summary]
\lean{Turing.StepBlock.summaryEmit_none}
\label{Turing.StepBlock.summaryEmit_none}
\leanok
\uses{Turing.StepBlock.summaryEmit}
\difficulty{3}
For every $b\in\mathrm{Option}(\mathrm{Bool})$, combining an empty summary with $b$
returns $b$.
\end{lemma}

\begin{definition}[Record width]
\lean{Turing.StepBlock.recordWidth}
\label{Turing.StepBlock.recordWidth}
\leanok
For each natural number $k$, define the record width to be $4k+4$.
\end{definition}

\begin{definition}[Record state space]
\lean{Turing.StepBlock.RecordState}
\label{Turing.StepBlock.RecordState}
\leanok
\uses{Turing.StepBlock.recordWidth}
For each natural number $k$, a record state is either absent or consists of a position
among $0,1,\ldots,4k+4$ together with a Boolean record of length $4k+4$.
\end{definition}

\begin{definition}[Table tape index]
\lean{Turing.StepBlock.tableTape}
\label{Turing.StepBlock.tableTape}
\leanok
For each natural number $k$, define the table tape index to be $k$, viewed as an index
among $k+2$ tape positions.
\end{definition}

\begin{definition}[Summary tape index]
\lean{Turing.StepBlock.summaryTape}
\label{Turing.StepBlock.summaryTape}
\leanok
For each natural number $k$, the summary tape is the position $k+1$ among the $k+2$ work
tapes, indexed from $0$.
\end{definition}

\begin{lemma}[Distinctness of an Original and Table Tape]
\lean{Turing.StepBlock.mirror_ne_table}
\label{Turing.StepBlock.mirror_ne_table}
\leanok
\uses{Turing.StepBlock.tableTape}
\difficulty{3}
For every natural number $k$ and every $i\in\{0,\ldots,k-1\}$, the tape indexed by $i$
among the original $k$ tapes is distinct from the additional table tape.
\end{lemma}

\begin{lemma}[Distinct summary and table tapes]
\lean{Turing.StepBlock.summary_ne_table}
\label{Turing.StepBlock.summary_ne_table}
\leanok
\uses{Turing.StepBlock.summaryTape, Turing.StepBlock.tableTape}
\difficulty{3}
For every nonnegative integer $k$, the summary tape among $k+2$ tape indices is
different from the table tape.
\end{lemma}

\begin{definition}[Encoded action payload]
\lean{Turing.StepBlock.payload}
\label{Turing.StepBlock.payload}
\leanok
\uses{Turing.Action, Turing.StepBlock.recordWidth, Turing.codeReadActionK}
For each natural number $k$ and each Boolean string of length $4k+4$, form the code-read
action for that string followed by a false bit, extract its action component, and use
the default action if no action is returned. The resulting action is for a machine with
$k$ work tapes, Boolean symbols, and a singleton state set.
\end{definition}

\begin{definition}[Applying an encoded payload]
\lean{Turing.StepBlock.applyPayload}
\label{Turing.StepBlock.applyPayload}
\leanok
\uses{Turing.Action, Turing.StepBlock.RecordState, Turing.StepBlock.summaryEmit}
Given an action on $k$ work tapes with Boolean symbols and a singleton state set, and an
optional Boolean summary symbol, form an action on $k+2$ work tapes with record states.
It copies the input-head movement and the original $k$ work-tape actions; on the first
additional tape it makes no write or movement, and on the second it stays in place and
writes the summary symbol obtained from the given symbol and the action's output. It
emits no output and enters the halting state.
\end{definition}

\begin{definition}[Record-reading Turing machine]
\lean{Turing.StepBlock.recordTM}
\label{Turing.StepBlock.recordTM}
\leanok
\uses{Turing.MultiTapeTM, Turing.StepBlock.RecordState, Turing.StepBlock.applyPayload, Turing.StepBlock.payload, Turing.StepBlock.recordWidth, Turing.StepBlock.summaryTape, Turing.StepBlock.tableTape, Turing.inputMoveAction, Turing.tapeAction}
For each natural number $k$, define a deterministic Turing machine with $k+2$ work
tapes, Boolean tape symbols, and states consisting either of a halting state or of a
cursor position together with a Boolean record of length $4k+4$. It starts with cursor
position $0$ and an all-false record. In a nonhalting state, if the cursor is below
$4k+4$, it reads the symbol under the designated table tape head (interpreting a blank
as false), stores that bit at the cursor position, advances the cursor, and moves that
tape head right. Once the cursor reaches the end of the record, it applies the encoded
payload using the symbol under the designated summary tape head. In the halting state it
makes no input-head movement.
\end{definition}

\begin{definition}[Reading a record into a configuration]
\lean{Turing.StepBlock.readCfg}
\label{Turing.StepBlock.readCfg}
\leanok
\uses{Turing.Cfg, Turing.StepBlock.RecordState, Turing.StepBlock.recordWidth, Turing.StepBlock.tableTape}
Let $k$ be a natural number, let $x$ be a Boolean input string, let $d$ be a
configuration on $k+2$ work tapes with record states, let $b$ be a Boolean record of
length $4k+4$, and let $r$ be a natural number with $r\le 4k+4$. Replace the state by
cursor position $r$ and the record that agrees with $b$ below position $r$ and is false
at every position at least $r$. Increase the head position on the designated table tape
by $r$; leave all other configuration data unchanged.
\end{definition}

\begin{lemma}[Correctness of the Encoded Action Payload]
\lean{Turing.StepBlock.payload_correct}
\label{Turing.StepBlock.payload_correct}
\leanok
\uses{Turing.Action, Turing.StepBlock.payload, Turing.StepBlock.recordWidth, Turing.actionBitsK, Turing.codeReadActionK_append, Turing.optBoolBits, Turing.optOptBoolBits, Turing.optStateBits, Turing.signBits}
\difficulty{5}
Let $k,n\in\mathbb N$, let $a$ be an action for a $k$-tape binary machine with states
indexed by $\{0,\ldots,n\}$, and let $bits$ assign a binary digit to each of the $4k+4$
positions in a record. If the word formed from these bits is the initial $4k+4$-bit
segment of the serialized action $a$, then decoding the record payload yields the same
input-head movement, work-tape actions, and output symbol as $a$, with no successor
state.
\end{lemma}

\begin{lemma}[Reading a Partial Record]
\lean{Turing.StepBlock.recordTM_read}
\label{Turing.StepBlock.recordTM_read}
\leanok
\uses{Turing.Cfg, Turing.Cfg.workTapeSymbols, Turing.MultiTapeTM.runFrom, Turing.MultiTapeTM.runFrom_succ_eq_step', Turing.StepBlock.RecordState, Turing.StepBlock.readCfg, Turing.StepBlock.recordTM, Turing.StepBlock.recordWidth, Turing.StepBlock.tableTape, Turing.step_of_tapeAction_move}
\difficulty{6}
Let $k\in\mathbb N$, let $x$ be a finite binary input, let $d$ be a configuration of the
$(k+2)$-work-tape record machine, and let $b$ be a Boolean record of width $4k+4$.
Suppose $d$ starts in the record machine's initial state and the record tape contains
$b$ at each of the next $4k+4$ cells from its head. For every $r\le 4k+4$, after $r$
steps the machine's configuration is exactly the configuration obtained from $d$ by
storing the record position $r$ and the record whose entries before position $r$ are
those of $b$ and whose remaining entries are false, and by advancing the record-tape
head by $r$.
\end{lemma}

\begin{lemma}[Execution of a Recorded Action]
\lean{Turing.StepBlock.recordTM_execute}
\label{Turing.StepBlock.recordTM_execute}
\leanok
\uses{Turing.Action, Turing.Action.apply, Turing.Cfg, Turing.Cfg.workTapeSymbols, Turing.MultiTapeTM.runFrom, Turing.MultiTapeTM.runFrom_succ_eq_step', Turing.MultiTapeTM.step, Turing.StepBlock.RecordState, Turing.StepBlock.applyPayload, Turing.StepBlock.payload_correct, Turing.StepBlock.readCfg, Turing.StepBlock.recordTM, Turing.StepBlock.recordTM_read, Turing.StepBlock.recordWidth, Turing.StepBlock.summaryTape, Turing.StepBlock.summary_ne_table, Turing.StepBlock.tableTape, Turing.actionBitsK}
\difficulty{4}
Let $k,n\in\mathbb N$, let $x$ be a finite binary input, let $d$ be a configuration of
the $(k+2)$-work-tape record machine, let $a$ be an action for a $k$-work-tape machine
with states indexed by $\{0,\ldots,n\}$, and let $b$ be a Boolean record of width
$4k+4$. Suppose $d$ starts in the record machine's initial state, the record tape
contains $b$ at each of the next $4k+4$ cells from its head, and $b$ is the
length-$4k+4$ prefix of the binary encoding of $a$. After $4k+5$ steps, the record
machine's configuration equals the one obtained by first setting the action's successor
state to halting, pairing that action with the current symbol on the summary tape to
form the payload action, and applying that payload action to the configuration whose
record state stores the full record and whose record-tape head is advanced by $4k+4$.
\end{lemma}

\begin{lemma}[Head-Movement Bounds During Recording]
\lean{Turing.StepBlock.recordTM_envelope}
\label{Turing.StepBlock.recordTM_envelope}
\leanok
\uses{Turing.Action, Turing.Cfg, Turing.Cfg.workTapeSymbols, Turing.MultiTapeTM.runFrom, Turing.StepBlock.RecordState, Turing.StepBlock.applyPayload, Turing.StepBlock.mirror_ne_table, Turing.StepBlock.readCfg, Turing.StepBlock.recordTM, Turing.StepBlock.recordTM_execute, Turing.StepBlock.recordTM_read, Turing.StepBlock.recordWidth, Turing.StepBlock.summaryTape, Turing.StepBlock.summary_ne_table, Turing.StepBlock.tableTape, Turing.actionBitsK, Turing.moveInputPos, Turing.workTapePos_apply_le}
\difficulty{6}
Let $k,n\in\mathbb N$, let $x$ be a finite binary input, let $d$ be a configuration of
the $(k+2)$-work-tape record machine, let $a$ be an action for a $k$-work-tape machine
with states indexed by $\{0,\ldots,n\}$, and let $b$ be a Boolean record of width
$4k+4$. Suppose $d$ starts in the record machine's initial state, the record tape
contains $b$ at each of the next $4k+4$ cells from its head, and $b$ is the
length-$4k+4$ prefix of the binary encoding of $a$. At every time $u\le 4k+5$, the first
$k$ work-tape heads have moved from their initial positions by the action's prescribed
displacements only when $u=4k+5$; the input head has moved according to $a$ only at that
final time and otherwise remains fixed. Every work-tape head has displacement in
absolute value at most $1$ on the first $k$ tapes and at most $4k+4$ on the final two
tapes, and the input-head displacement is at most $1$. The record-tape head is at its
initial position plus $\min(u,4k+4)$, and the summary-tape head remains at its initial
position.
\end{lemma}

\begin{definition}[Finite record-reading machine]
\lean{Turing.StepBlock.recordFinTM}
\label{Turing.StepBlock.recordFinTM}
\leanok
\uses{Turing.FinTM, Turing.StepBlock.RecordState, Turing.StepBlock.recordTM}
For each natural number $k$, define a finite Turing machine over the Boolean alphabet
with $k+2$ work tapes and record states consisting of a halting state or a cursor
position and a Boolean record of length $4k+4$. Its underlying machine is the
record-reading machine, and its state set is finite with decidable equality.
\end{definition}

\begin{lemma}[Effect of the Record Machine]
\lean{Turing.StepBlock.recordTM_effect}
\label{Turing.StepBlock.recordTM_effect}
\leanok
\uses{Turing.Action, Turing.Action.apply, Turing.Cfg, Turing.Cfg.workTapeSymbols, Turing.MultiTapeTM.runFrom, Turing.StepBlock.RecordState, Turing.StepBlock.applyPayload, Turing.StepBlock.mirror_ne_table, Turing.StepBlock.outputSummary, Turing.StepBlock.outputSummary_append, Turing.StepBlock.readCfg, Turing.StepBlock.recordTM, Turing.StepBlock.recordTM_envelope, Turing.StepBlock.recordTM_execute, Turing.StepBlock.recordTM_read, Turing.StepBlock.recordWidth, Turing.StepBlock.summaryTape, Turing.StepBlock.tableTape, Turing.actionBitsK, Turing.codeActionK_length, Turing.moveInputPos}
\difficulty{5}
Let $k,n\in\mathbb N$, let $x$ be a finite binary input, and let $d$ be a configuration
of a $(k+2)$-work-tape machine with record states, while $c$ is a configuration of a
$k$-work-tape machine and $a$ is one of its actions. Let $b$ be a Boolean record of
width $4k+4$. Suppose $d$ starts in the record machine's initial state, the record tape
contains $b$ at each of the next $4k+4$ cells from its head, and $b$ is the prefix of
length $4k+4$ of the binary encoding of $a$. Suppose also that $d$ and $c$ have the same
input-head position, their first $k$ work tapes and head positions agree, and the
summary tape of $d$ currently reads the output summary of $c$. After $4k+5$ steps, the
record machine is halted; its input-head position, first $k$ work tapes and their head
positions, and summary-tape symbol agree with those obtained by applying $a$ to $c$,
while its output list is unchanged from $d$. It does not halt at any earlier time; the
encoded action has length at least $4k+5$; and at every time $u\le 4k+5$, the original
work-tape heads have moved by their prescribed action displacements only at the final
time, every work-tape head has moved by at most $1$ on the first $k$ tapes and at most
$4k+4$ on the other two, the input head has moved by at most $1$, the record-tape head
has advanced by $\min(u,4k+4)$, and the summary-tape head has not moved.
\end{lemma}

\begin{lemma}[Completeness of work-read patterns]
\lean{Turing.StepBlock.mem_workReadsK}
\label{Turing.StepBlock.mem_workReadsK}
\leanok
\uses{Turing.workReadsK}
\difficulty{4}
For every natural number $k$, every function $w : 	ext{Fin}(k) 	o
ext{Option}(	ext{Bool})$ occurs in the list obtained by recursively listing all
length-$k$ sequences whose entries are either absent, false, or true.
\end{lemma}

\begin{lemma}[Number of work-tape read patterns]
\lean{Turing.StepBlock.length_workReadsK}
\label{Turing.StepBlock.length_workReadsK}
\leanok
\uses{Turing.workReadsK}
\difficulty{4}
For every natural number $k$, the list of possible read patterns on $k$ binary work
tapes, with each scanned cell either blank or containing one of the two bits, has length
$3^k$.
\end{lemma}

\begin{definition}[Rank of a work-tape reading pattern]
\lean{Turing.StepBlock.workRank}
\label{Turing.StepBlock.workRank}
\leanok
\uses{Turing.StepBlock.length_workReadsK, Turing.StepBlock.mem_workReadsK, Turing.workReadsK}
For each natural number $k$, assign a work-tape reading pattern—one optional Boolean
symbol for each of the $k$ tapes—the position of that pattern in the ordered list of all
such patterns. This rank lies between $0$ and $3^k-1$.
\end{definition}

\begin{definition}[Rank of an input symbol]
\lean{Turing.StepBlock.inputRank}
\label{Turing.StepBlock.inputRank}
\leanok
Assign rank $0$ to a blank input cell, rank $1$ to a cell containing false, and rank $2$
to a cell containing true.
\end{definition}

\begin{definition}[Rank of an input and work-tape reading]
\lean{Turing.StepBlock.readRank}
\label{Turing.StepBlock.readRank}
\leanok
\uses{Turing.StepBlock.inputRank, Turing.StepBlock.workRank}
For a machine reading an input symbol that is either a Boolean value or blank, together
with Boolean-or-blank symbols on each of its $k$ work tapes, the read rank is the
element of ${0,\ldots,3^{k+1}-1}$ obtained by encoding the input symbol in base $3$ as
the leading digit and the work-tape symbols as the remaining rank.
\end{definition}

\begin{definition}[Successor-processing phases]
\lean{Turing.StepBlock.SuccPhase}
\label{Turing.StepBlock.SuccPhase}
\leanok
The successor-processing phase is one of three forms: a tag phase, a unary phase, or a
completed phase carrying a Boolean value indicating whether the successor is live.
\end{definition}

\begin{definition}[Successor-field transition machine]
\lean{Turing.StepBlock.succFieldTM}
\label{Turing.StepBlock.succFieldTM}
\leanok
\uses{Turing.MultiTapeTM, Turing.StepBlock.SuccPhase, Turing.inputMoveAction, Turing.tapeAction}
For a natural number $m$, two work-tape indices
$\mathrm{src},\mathrm{dst}\in\{0,\ldots,m-1\}$, and a Boolean choice $c$, define an
$m$-work-tape machine with Boolean tape alphabet and the three successor-processing
phases. It starts in the tag phase. In that phase it moves the source head one cell
right and enters the unary phase if the scanned source symbol is $1$, and otherwise
completes with the live flag false. In the unary phase it moves the source head right,
and, when $c$ is true, writes the scanned source symbol on the destination tape; it
remains in the unary phase if that symbol is $1$, and otherwise completes with the live
flag true. In a completed phase it stays put on all tapes and preserves that phase and
its flag.
\end{definition}

\begin{definition}[Successor-field configuration update]
\lean{Turing.StepBlock.succFieldCfg}
\label{Turing.StepBlock.succFieldCfg}
\leanok
\uses{Turing.Cfg, Turing.FinTM.bufferTape, Turing.StepBlock.SuccPhase}
Given a natural number $m$, a Boolean input string, source and destination work-tape
indices, a Boolean choice $c$, an initial configuration $d$, a Boolean string $w$, a
natural number $r$, and a successor-processing phase $q$, define a new configuration
with the same input-head position and output as $d$ and state $q$. If $c$ is true,
replace the destination tape's cells in the interval from its old head position
(inclusive) to that position plus $r$ (exclusive) by the corresponding symbols of $w$,
interpreting positions in $w$ beyond its length as blank; leave every other work-tape
cell unchanged. Advance the source head by $r$, and also advance the destination head by
$r$ exactly when $c$ is true and the destination and source indices coincide; leave all
other work-head positions unchanged.
\end{definition}

\begin{lemma}[Unary copying step]
\lean{Turing.StepBlock.succFieldTM_unary_step}
\label{Turing.StepBlock.succFieldTM_unary_step}
\leanok
\uses{Turing.Action, Turing.Action.apply, Turing.Cfg, Turing.Cfg.workTapeSymbols, Turing.MultiTapeTM.step, Turing.StepBlock.SuccPhase, Turing.StepBlock.succFieldCfg, Turing.StepBlock.succFieldTM, Turing.catalog_write_take}
\difficulty{5}
Let $m$ be a nonnegative integer, let $x$ be a Boolean input word, and choose distinct
work-tape indices $s,t\in\{0,\ldots,m-1\}$. Let $c$ indicate whether copying is enabled,
and let $d$ be a configuration with $m$ Boolean work tapes. For a Boolean word $w$ and
an integer $r<|w|$, suppose tape $s$ in $d$ contains the symbol $w_r$ at the cell $r$
positions to the right of its head. Form a configuration from $d$ by setting its phase
to unary, advancing the heads on tape $s$ and, when copying is enabled, tape $t$ by $r$,
and placing the first $r$ symbols of $w$ on tape $t$ starting at its original head
position when copying is enabled. One step of the machine so defined produces the
configuration formed in the same way with $r+1$ in place of $r$, and with phase unary if
$w_r$ is true and done otherwise.
\end{lemma}

\begin{lemma}[Run through a unary state field]
\lean{Turing.StepBlock.succFieldTM_unary_run}
\label{Turing.StepBlock.succFieldTM_unary_run}
\leanok
\uses{Turing.Cfg, Turing.MultiTapeTM.runFrom, Turing.MultiTapeTM.runFrom_succ_eq_step', Turing.StepBlock.SuccPhase, Turing.StepBlock.succFieldCfg, Turing.StepBlock.succFieldTM, Turing.StepBlock.succFieldTM_unary_step}
\difficulty{5}
Let $m$ be a natural number, let $x$ be a finite Boolean input, let $src,dst$ be
distinct work-tape indices, let $copy$ be a Boolean, and let $d$ be a configuration of
the successor-field machine. Suppose $d$ is in the unary phase and the source tape,
starting at its head position, contains $s$ ones followed by a zero. For any $r\le s+1$,
after $r$ steps the machine is in the configuration with the first $r$ bits of that word
copied to the destination tape exactly when copying is enabled; its state is unary if
$r\le s$, and done with the live flag true if $r=s+1$.
\end{lemma}

\begin{definition}[Configuration after the tag phase]
\lean{Turing.StepBlock.afterTag}
\label{Turing.StepBlock.afterTag}
\leanok
\uses{Turing.Cfg, Turing.StepBlock.SuccPhase}
Given a natural number $m$, a Boolean input string, a source work-tape index, and a
configuration, define the configuration with all fields unchanged except that its state
is the unary phase and the source work-tape head is advanced one cell to the right.
\end{definition}

\begin{lemma}[Successor-field run at each intermediate time]
\lean{Turing.StepBlock.succFieldTM_trace}
\label{Turing.StepBlock.succFieldTM_trace}
\leanok
\uses{Turing.Cfg, Turing.Cfg.workTapeSymbols, Turing.MultiTapeTM.runFrom, Turing.MultiTapeTM.runFrom_succ_eq_step, Turing.MultiTapeTM.step, Turing.StepBlock.SuccPhase, Turing.StepBlock.afterTag, Turing.StepBlock.succFieldCfg, Turing.StepBlock.succFieldTM, Turing.StepBlock.succFieldTM_unary_run, Turing.optStateBits, Turing.step_of_tapeAction_move, Turing.unaryFin}
\difficulty{6}
Let $m,n$ be natural numbers, let $x$ be a finite Boolean input, let $src,dst$ be
distinct work-tape indices, let $copy$ be a Boolean, and let $d$ be a configuration of
the successor-field machine. Let $s$ be absent or a state index less than $n$. Suppose
$d$ is in the tag state and the source tape, starting at its head position, contains the
optional-state encoding of $s$. For every time $t$ no greater than the encoding length,
the configuration after $t$ steps is $d$ if $t=0$. Otherwise it is the configuration
obtained after the tag transition, with the prefix of length $t-1$ of the unary state
encoding copied to the destination tape exactly when copying is enabled, and with state
unary when $t$ is strictly less than the encoding length and done (carrying whether $s$
is present) otherwise.
\end{lemma}

\begin{lemma}[Completed successor-field run]
\lean{Turing.StepBlock.succFieldTM_run}
\label{Turing.StepBlock.succFieldTM_run}
\leanok
\uses{Turing.Cfg, Turing.MultiTapeTM.runFrom, Turing.StepBlock.SuccPhase, Turing.StepBlock.afterTag, Turing.StepBlock.succFieldCfg, Turing.StepBlock.succFieldTM, Turing.StepBlock.succFieldTM_trace, Turing.optStateBits, Turing.unaryFin}
\difficulty{4}
Let $m,n$ be natural numbers, let $x$ be a finite Boolean input, let $src,dst$ be
distinct work-tape indices, let $copy$ be a Boolean, and let $d$ be a configuration of
the successor-field machine. Let $s$ be absent or a state index less than $n$. Suppose
$d$ is in the tag state and the source tape, starting at its head position, contains the
optional-state encoding of $s$. After as many steps as the length of that encoding, the
machine reaches the configuration obtained by advancing from the tag state to the unary
phase, with the unary encoding of $s$ copied to the destination tape exactly when $s$ is
present and copying is enabled, and with the state set to done carrying whether $s$ is
present.
\end{lemma}

\begin{lemma}[No early completion while scanning a successor field]
\lean{Turing.StepBlock.succFieldTM_cut}
\label{Turing.StepBlock.succFieldTM_cut}
\leanok
\uses{Turing.Cfg, Turing.MultiTapeTM.runFrom, Turing.StepBlock.SuccPhase, Turing.StepBlock.succFieldCfg, Turing.StepBlock.succFieldTM, Turing.StepBlock.succFieldTM_trace, Turing.optStateBits}
\difficulty{4}
Let $m,n$ be natural numbers, let $x$ be a finite Boolean input, let $src,dst$ be
distinct work-tape indices, let $copy$ be a Boolean, and let $d$ be a configuration of
the successor-field machine. Let $s$ be either absent or a state index less than $n$.
Suppose $d$ is in the tag state and the source tape, starting at its head position,
contains the complete optional-state encoding of $s$. At every time $t$ strictly less
than the length of that encoding, the machine's state is not any done state.
\end{lemma}

\begin{lemma}[Head positions during successor-field processing]
\lean{Turing.StepBlock.succFieldTM_heads}
\label{Turing.StepBlock.succFieldTM_heads}
\leanok
\uses{Turing.Cfg, Turing.MultiTapeTM.runFrom, Turing.StepBlock.SuccPhase, Turing.StepBlock.afterTag, Turing.StepBlock.succFieldCfg, Turing.StepBlock.succFieldTM, Turing.StepBlock.succFieldTM_trace, Turing.optStateBits}
\difficulty{4}
Let $m,n$ be natural numbers, let $x$ be a finite Boolean input, let $src,dst$ be
distinct work-tape indices, let $copy$ be a Boolean, and let $d$ be a configuration of
the successor-field machine. Let $s$ be absent or a state index less than $n$. Suppose
$d$ is in the tag state and the source tape, starting at its head position, contains the
optional-state encoding of $s$. At each time $t$ no greater than the encoding length,
the input-head position and output agree with those of $d$. Every work-tape head
advances by $t$ on the source tape, by $t-1$ on the destination tape if copying is
enabled, and by zero on all other tapes.
\end{lemma}

\begin{definition}[Phase-dependent state type]
\lean{Turing.StepBlock.SkipPh}
\label{Turing.StepBlock.SkipPh}
\leanok
\uses{Turing.StepBlock.RecordState, Turing.StepBlock.SuccPhase}
For each natural number $k$, define a state type depending on a Boolean parameter: when
the parameter is false, it is the record-state type for width $4k+4$; when true, it is
the successor-processing phase type.
\end{definition}

\begin{definition}[Machine for skipping a record part]
\lean{Turing.StepBlock.skipPart}
\label{Turing.StepBlock.skipPart}
\leanok
\uses{Turing.MultiTapeTM, Turing.StepBlock.SkipPh, Turing.StepBlock.recordTM, Turing.StepBlock.succFieldTM, Turing.StepBlock.summaryTape, Turing.StepBlock.tableTape}
For each natural number $k$ and Boolean parameter $b$, define a machine with $k+2$ work
tapes, Boolean tape alphabet, and the state type selected by $b$. When $b$ is false,
this machine is the record-processing machine of width $4k+4$; when $b$ is true, it is
the successor-field machine with the table tape as source, the summary tape as
destination, and copying disabled.
\end{definition}

\begin{definition}[Reduction at the record boundary]
\lean{Turing.StepBlock.skipRed}
\label{Turing.StepBlock.skipRed}
\leanok
\uses{Turing.StepBlock.SkipPh, Turing.StepBlock.recordWidth}
For each natural number $k$, define a partial state reduction on the phase-dependent
state type. In the false-parameter case, a record state carrying a position equal to the
record width $4k+4$ is sent to the true-parameter state type with tag phase; all other
states, including the absent record state, are sent to no value. In the true-parameter
case, every state is sent to no value.
\end{definition}

\begin{definition}[Record-and-skip machine]
\lean{Turing.StepBlock.recordSkipTM}
\label{Turing.StepBlock.recordSkipTM}
\leanok
\uses{Turing.MultiTapeTM, Turing.StepBlock.SkipPh, Turing.StepBlock.recordTM, Turing.StepBlock.skipPart, Turing.StepBlock.skipRed, Turing.phaseHostTM}
For each natural number $k$, define a deterministic Turing machine with $k+2$ work
tapes, Boolean tape symbols, and states tagged by a Boolean phase. In the false phase
its states are those of the record machine; when that machine reaches position $4k+4$,
the phase changes to true and the machine enters the successor-field routine in its
initial tag state. This is the phase-hosted machine built from the record machine and
the skip reduction.
\end{definition}

\begin{lemma}[Fixed-width recording and skip run]
\lean{Turing.StepBlock.recordSkipTM_fixed}
\label{Turing.StepBlock.recordSkipTM_fixed}
\leanok
\uses{Turing.Cfg, Turing.Cfg.mapState, Turing.MultiTapeTM.runFrom, Turing.StepBlock.RecordState, Turing.StepBlock.SkipPh, Turing.StepBlock.readCfg, Turing.StepBlock.recordSkipTM, Turing.StepBlock.recordTM, Turing.StepBlock.recordTM_read, Turing.StepBlock.recordWidth, Turing.StepBlock.skipPart, Turing.StepBlock.skipRed, Turing.StepBlock.tableTape, Turing.phaseHostTM_run}
\difficulty{5}
Let $k\in\mathbb{N}$, let $x$ be a finite Boolean word, and let $d$ be a configuration
with $k+2$ work tapes, Boolean tape symbols, and record states. Let $\mathrm{bits}$
assign a bit to each of the $4k+4$ record positions. Suppose $d$ is in the initial state
of the recording machine and, for every record position $i$, the table tape at the
current table-head position plus $i$ contains $\mathrm{bits}(i)$. For each $u\le 4k+4$,
running the record-and-skip machine for $u$ steps from the state-renamed configuration
gives the state-renamed configuration that has read the first $u$ record bits (with the
prescribed read configuration); after $4k+5$ steps, it gives the read configuration at
full width, with its state renamed to the successor phase tagged as true.
\end{lemma}

\begin{lemma}[Run in the successor phase]
\lean{Turing.StepBlock.recordSkipTM_tail}
\label{Turing.StepBlock.recordSkipTM_tail}
\leanok
\uses{Turing.Cfg, Turing.Cfg.mapState, Turing.MultiTapeTM.runFrom, Turing.MultiTapeTM.runFrom_mapState_of_agreeOn, Turing.StepBlock.SkipPh, Turing.StepBlock.SuccPhase, Turing.StepBlock.recordSkipTM, Turing.StepBlock.succFieldTM, Turing.StepBlock.summaryTape, Turing.StepBlock.tableTape}
\difficulty{1}
Let $k\in\mathbb{N}$, let $x$ be a finite Boolean word, and let $d$ be a configuration
with $k+2$ work tapes, Boolean symbols, and successor-phase states. For every natural
number $t$, running the record-and-skip machine for $t$ steps from the state-renamed
configuration equals running the successor-field machine for $t$ steps from $d$ and then
renaming its state.
\end{lemma}

\begin{lemma}[Fixed-length action prefix]
\lean{Turing.StepBlock.actionBitsK_split}
\label{Turing.StepBlock.actionBitsK_split}
\leanok
\uses{Turing.Action, Turing.StepBlock.recordWidth, Turing.actionBitsK, Turing.optBoolBits, Turing.optOptBoolBits, Turing.optStateBits, Turing.signBits}
\difficulty{4}
Let $k,n$ be natural numbers, and let $a$ be an action of a $k$-tape Turing machine with
Boolean work-tape symbols and successor states in $	ext{Fin}(n+1)$. Its action-bit
encoding is the concatenation of a Boolean list of length $4k+4$ and the encoding of its
optional successor state.
\end{lemma}

\begin{lemma}[Reads on concatenated words]
\lean{Turing.StepBlock.reads_append}
\label{Turing.StepBlock.reads_append}
\leanok
\difficulty{3}
Let $c:\mathbb{Z}\to\mathrm{Option}(\mathrm{Bool})$, let $p\in\mathbb{Z}$, and let $a,b$
be finite Boolean words. Suppose that for every natural number $r$ less than the length
of $a$ concatenated with $b$, the cell at position $p+r$ contains the $r$-th bit of that
concatenation. Then every position $p+r$ for $r$ within $a$ contains the corresponding
bit of $a$, and every position $p+|a|+r$ for $r$ within $b$ contains the corresponding
bit of $b$.
\end{lemma}

\begin{lemma}[Action record and table reads]
\lean{Turing.StepBlock.actionBitsK_reads}
\label{Turing.StepBlock.actionBitsK_reads}
\leanok
\uses{Turing.Action, Turing.Cfg, Turing.StepBlock.actionBitsK_split, Turing.StepBlock.reads_append, Turing.StepBlock.recordWidth, Turing.StepBlock.tableTape, Turing.actionBitsK, Turing.optStateBits}
\difficulty{4}
Let $k,n$ be natural numbers, let $S$ be a type, let $x$ be a finite input list of
symbols, let $d$ be a configuration with $k+2$ work tapes, Boolean tape symbols, and
state type $S$, and let $a$ be a $k$-tape action with Boolean work-tape symbols and
successor states in $	ext{Fin}(n+1)$. Suppose the extra work tape, starting at its head
position, contains the serialized action $a$ with no blank among its bits. Then there is
a Boolean record of length $4k+4$ such that the action serialization is this record
followed by the successor-state serialization, the total length is $4k+4$ plus the
length of that state serialization, the record bits are read from the corresponding
consecutive cells of the extra tape, and the remaining cells read the corresponding
successor-state bits.
\end{lemma}

\begin{lemma}[Completed recording and skip run]
\lean{Turing.StepBlock.recordSkipTM_run}
\label{Turing.StepBlock.recordSkipTM_run}
\leanok
\uses{Turing.Action, Turing.Cfg, Turing.Cfg.mapState, Turing.MultiTapeTM.runFrom, Turing.MultiTapeTM.runFrom_add, Turing.StepBlock.RecordState, Turing.StepBlock.SkipPh, Turing.StepBlock.SuccPhase, Turing.StepBlock.actionBitsK_reads, Turing.StepBlock.afterTag, Turing.StepBlock.readCfg, Turing.StepBlock.recordSkipTM, Turing.StepBlock.recordSkipTM_fixed, Turing.StepBlock.recordSkipTM_tail, Turing.StepBlock.recordTM, Turing.StepBlock.recordWidth, Turing.StepBlock.succFieldCfg, Turing.StepBlock.succFieldTM_run, Turing.StepBlock.summary_ne_table, Turing.StepBlock.tableTape, Turing.actionBitsK, Turing.optStateBits, Turing.unaryFin}
\difficulty{5}
Let $k,n\in\mathbb{N}$, let $x$ be a finite Boolean word, let $d$ be a configuration
with $k+2$ work tapes, Boolean symbols, and record states, and let $a$ be an action
whose successor state is in $\{0,\ldots,n\}$ or absent. Suppose $d$ has the initial
recording-machine state and its table tape contains, in order from its current head
position, the bits encoding $a$. After one more step than the length of that encoding,
the record-and-skip machine has the same input-head position, work-tape contents,
output, and all work-tape head positions as $d$, except that the table-tape head has
advanced by the encoding length. Its state is the successor phase marked done, with the
Boolean value indicating whether $a$ has a successor state.
\end{lemma}

\begin{lemma}[Path of the Recording and Skip Machine]
\lean{Turing.StepBlock.recordSkipTM_path}
\label{Turing.StepBlock.recordSkipTM_path}
\leanok
\uses{Turing.Cfg, Turing.Cfg.mapState, Turing.MultiTapeTM.runFrom, Turing.MultiTapeTM.runFrom_add, Turing.StepBlock.RecordState, Turing.StepBlock.SkipPh, Turing.StepBlock.SuccPhase, Turing.StepBlock.readCfg, Turing.StepBlock.recordSkipTM, Turing.StepBlock.recordSkipTM_fixed, Turing.StepBlock.recordSkipTM_tail, Turing.StepBlock.recordTM, Turing.StepBlock.recordWidth, Turing.StepBlock.succFieldTM, Turing.StepBlock.succFieldTM_cut, Turing.StepBlock.succFieldTM_heads, Turing.StepBlock.succFieldTM_trace, Turing.StepBlock.summaryTape, Turing.StepBlock.summary_ne_table, Turing.StepBlock.tableTape, Turing.optStateBits}
\difficulty{6}
Let $k,n$ be natural numbers, let $x$ be a finite Boolean word, and let $d$ be a
configuration of the recording machine with $k+2$ work tapes and Boolean symbols. Let
$b$ assign a bit to each of the $4k+4$ record positions, and let $s$ be either an index
in $\\{0,\ldots,n-1\}$ or absent. Suppose $d$ is in the recording machine's initial
state, its table tape contains $b(i)$ at the cell offset by $i$ from its head for each
record position $i$, and the cells immediately following the record contain the bit
serialization of $s$. For every natural number $u\le
4k+5+|\operatorname{optStateBits}(s)|$, after $u$ steps of the recording-and-skip
machine the input-head position, all work-tape contents, and output are unchanged; each
work-tape head has moved by $\min(u,4k+4)+\max(u-(4k+5),0)$ on the table tape and by
zero on every other tape. If $u<4k+5+|\operatorname{optStateBits}(s)|$, the machine is
not in the tagged successor phase with either done value.
\end{lemma}

\begin{lemma}[Recording-machine run specification]
\lean{Turing.StepBlock.recordSkipTM_spec}
\label{Turing.StepBlock.recordSkipTM_spec}
\leanok
\uses{Turing.Action, Turing.Cfg, Turing.Cfg.mapState, Turing.MultiTapeTM.runFrom, Turing.StepBlock.RecordState, Turing.StepBlock.actionBitsK_reads, Turing.StepBlock.recordSkipTM, Turing.StepBlock.recordSkipTM_path, Turing.StepBlock.recordTM, Turing.StepBlock.recordWidth, Turing.StepBlock.tableTape, Turing.actionBitsK}
\difficulty{3}
Let $k,n\in\mathbb{N}$, let $x$ be a finite Boolean word, let $d$ be a configuration
with $k+2$ work tapes, Boolean symbols, and record states, and let $a$ be an action
whose successor state is in $\{0,\ldots,n\}$ or absent. Suppose $d$ has the initial
recording-machine state and its table tape contains the bits encoding $a$ in order from
its current head position. For every $u$ less than the encoding length plus one, the
machine's state after $u$ steps is not a tagged done state. For every $u$ up to and
including the encoding length plus one, its input-head position, work-tape contents, and
output equal their initial values; the head on the table tape is shifted by
$\min(u,4k+4)+\max(u-(4k+5),0)$, and every other work-tape head is unchanged.
\end{lemma}

\begin{definition}[Seek control states]
\lean{Turing.StepBlock.SeekCtrl}
\label{Turing.StepBlock.SeekCtrl}
\leanok
For each natural number $k$, define the seek-control state type to have states
$	ext{observe}$ and $	ext{done}$, states $	ext{advance}(r)$ and $	ext{state}(r)$ for
each rank $r$ with $0 \le r < 3^{k+1}$, and states $	ext{request}(r,s,b)$ for ranks
$r,s$ in that range and a Boolean $b$.
\end{definition}

\begin{definition}[Seek router machine]
\lean{Turing.StepBlock.seekRouterTM}
\label{Turing.StepBlock.seekRouterTM}
\leanok
\uses{Turing.MultiTapeTM, Turing.StepBlock.SeekCtrl, Turing.StepBlock.readRank, Turing.StepBlock.summaryTape, Turing.inputMoveAction, Turing.tapeAction}
For each natural number $k$, define a deterministic Turing machine with $k+2$ work
tapes, Boolean symbols, and seek-control states, initially in the observe state. It
computes the rank of the current input symbol together with the symbols under the first
$k$ work-tape heads; it then advances the summary-tape head and examines its scanned
symbol. If that symbol is true, it advances the summary-tape head again and requests the
computed rank with the remaining rank set to $3^{k+1}-1$ and the Boolean flag false. If
it is not true, it finishes when the computed rank is zero, and otherwise requests the
rank one smaller with the flag true. In request and done states it leaves the input head
stationary and keeps the same control state.
\end{definition}

\begin{definition}[Next seek state]
\lean{Turing.StepBlock.seekNext}
\label{Turing.StepBlock.seekNext}
\leanok
\uses{Turing.StepBlock.SeekCtrl}
For a natural number $k$, ranks $r,s$ with $0 \le r,s < 3^{k+1}$, and a Boolean $b$,
define the next seek state as follows: if $s=0$, it is done when $b$ is true and is the
state with rank $r$ when $b$ is false; if $s>0$, it is the request state with ranks $r$
and $s-1$ and flag $b$.
\end{definition}

\begin{definition}[Seek keys]
\lean{Turing.StepBlock.SeekKey}
\label{Turing.StepBlock.SeekKey}
\leanok
For each natural number $k$, define a seek key to be either absent or a triple
consisting of two ranks $r,s$ with $0 \le r,s < 3^{k+1}$ and a Boolean.
\end{definition}

\begin{definition}[Seek phase types]
\lean{Turing.StepBlock.SeekPh}
\label{Turing.StepBlock.SeekPh}
\leanok
\uses{Turing.StepBlock.SeekCtrl, Turing.StepBlock.SeekKey, Turing.StepBlock.SkipPh}
For each natural number $k$, define a type depending on a seek key: for an absent key it
is the seek-control state type, and for a present key it is the type of pairs consisting
of a Boolean phase and a state in the corresponding skip-phase type.
\end{definition}

\begin{definition}[Seek machine components]
\lean{Turing.StepBlock.seekPart}
\label{Turing.StepBlock.seekPart}
\leanok
\uses{Turing.MultiTapeTM, Turing.StepBlock.SeekKey, Turing.StepBlock.SeekPh, Turing.StepBlock.recordSkipTM, Turing.StepBlock.seekRouterTM}
For each natural number $k$, define a machine depending on a seek key. For an absent
key, it is the seek-router machine; for a present key, it is the record-and-skip
machine.
\end{definition}

\begin{definition}[Seek phase reduction]
\lean{Turing.StepBlock.seekRed}
\label{Turing.StepBlock.seekRed}
\leanok
\uses{Turing.StepBlock.SeekKey, Turing.StepBlock.SeekPh, Turing.StepBlock.recordSkipTM, Turing.StepBlock.seekNext}
For each natural number $k$, define a reduction on seek keys and their corresponding
phase states. From an absent key and a request state with ranks $r,s$ and flag $b$, it
returns the key $(r,s,b)$ together with the initial state of the record-and-skip
machine. From a present key $(r,s,b)$ and a true-phase done state, it returns the absent
key together with the next seek state determined by $r,s,b$. In every other case it
returns no reduction.
\end{definition}

\begin{definition}[Seek machine construction]
\lean{Turing.StepBlock.seekTM}
\label{Turing.StepBlock.seekTM}
\leanok
\uses{Turing.MultiTapeTM, Turing.StepBlock.SeekPh, Turing.StepBlock.seekPart, Turing.StepBlock.seekRed, Turing.phaseHostTM}
For each natural number $k$, the seek machine is the $(k+2)$-work-tape deterministic
Turing machine over the Boolean alphabet whose state is a pair consisting of an optional
seek key and a phase indexed by that key. It is constructed from the seek component,
starts in the state with no key and the observe control, and uses the seek reduction.
\end{definition}

\begin{definition}[Seek control configuration]
\lean{Turing.StepBlock.seekCfg}
\label{Turing.StepBlock.seekCfg}
\leanok
\uses{Turing.Cfg, Turing.Cfg.mapState, Turing.StepBlock.RecordState, Turing.StepBlock.SeekCtrl, Turing.StepBlock.SeekPh}
Let $k$ be a natural number, let $x$ be a finite Boolean input string, and let $d$ be a
configuration with $k+2$ work tapes, Boolean tape symbols, and record states. For any
seek control $q$, the associated seek-control configuration has the same tape data, head
positions, and output as $d$, and its state is the pair consisting of no seek key and
$q$.
\end{definition}

\begin{lemma}[Request processing in the seek machine]
\lean{Turing.StepBlock.seekTM_request}
\label{Turing.StepBlock.seekTM_request}
\leanok
\uses{Turing.Action, Turing.Cfg, Turing.Cfg.mapState, Turing.MultiTapeTM.runFrom, Turing.MultiTapeTM.runFrom_succ_eq_step, Turing.MultiTapeTM.step, Turing.StepBlock.RecordState, Turing.StepBlock.SeekKey, Turing.StepBlock.SkipPh, Turing.StepBlock.recordSkipTM, Turing.StepBlock.recordSkipTM_run, Turing.StepBlock.recordSkipTM_spec, Turing.StepBlock.recordTM, Turing.StepBlock.seekCfg, Turing.StepBlock.seekNext, Turing.StepBlock.seekPart, Turing.StepBlock.seekRed, Turing.StepBlock.seekTM, Turing.StepBlock.tableTape, Turing.actionBitsK, Turing.phaseHostTM, Turing.phaseHostTM_run, Turing.step_of_inputMoveAction}
\difficulty{6}
Let $k,n$ be natural numbers, let $x$ be a finite Boolean input, let $d$ be a
configuration of the record machine, and let $a$ be an action for $k$ work tapes and
$n+1$ states. Let $rank$ and $rest$ be indices less than $3^{k+1}$, and let $last$ be a
Boolean. Suppose the state of $d$ is the record machine's initial state and, at every
position from the table head through the length of the action's Boolean encoding, the
table tape contains that encoding. After the seek machine runs for the encoding length
plus three steps from the request configuration, it reaches the configuration with the
table head advanced by that length and control set to the successor of the request
specified by $rank$, $rest$, and $last$. Before that time its state is never the
terminal control state; throughout those steps, the input-head position, tape contents,
output, and all work-tape head positions other than the table head remain unchanged.
\end{lemma}

\begin{definition}[Cost of a list of actions]
\lean{Turing.StepBlock.seekCost}
\label{Turing.StepBlock.seekCost}
\leanok
\uses{Turing.Action, Turing.actionBitsK}
Let $k$ and $n$ be natural numbers, and let a finite list consist of actions for a
$k$-tape machine with Boolean symbols and states in the set of $n+1$ states. Its seek
cost is the length of the concatenation of the action-bit encodings of the actions, plus
three times the number of actions.
\end{definition}

\begin{lemma}[Search through a list of records]
\lean{Turing.StepBlock.seekTM_many}
\label{Turing.StepBlock.seekTM_many}
\leanok
\uses{Turing.Action, Turing.Cfg, Turing.Cfg.mapState, Turing.MultiTapeTM.runFrom, Turing.MultiTapeTM.runFrom_add, Turing.StepBlock.RecordState, Turing.StepBlock.reads_append, Turing.StepBlock.recordTM, Turing.StepBlock.seekCfg, Turing.StepBlock.seekCost, Turing.StepBlock.seekNext, Turing.StepBlock.seekTM, Turing.StepBlock.seekTM_request, Turing.StepBlock.tableTape, Turing.actionBitsK}
\difficulty{6}
Let $k,n$ be nonnegative integers, let $x$ be a finite binary input, and let $A$ be a
nonempty list of transition actions with at most $3^{k+1}$ entries. Let $d$ be a
configuration whose state is the recording machine's initial state, and suppose its
table tape stores the bit encoding of $A$ consecutively from the head. For any control
rank $r$ and Boolean flag $b$, start the search machine in the request state with rank
$r$, remaining count $|A|-1$, and flag $b$. After $\operatorname{cost}(A)$ steps, its
remaining count has been processed and it is in the done state if $b$ is true, and
otherwise in the state carrying rank $r$. Its table head has advanced by the encoding
length of $A$, and all other tape contents, the input-head position, and the output are
unchanged; every other work-tape head remains fixed. It does not reach the done state at
an earlier time.
\end{lemma}

\begin{lemma}[Seek machine state transition]
\lean{Turing.StepBlock.seekTM_stateStep}
\label{Turing.StepBlock.seekTM_stateStep}
\leanok
\uses{Turing.Action.mapState, Turing.Cfg, Turing.Cfg.mapState, Turing.Cfg.workTapeSymbols, Turing.MultiTapeTM.step, Turing.StepBlock.RecordState, Turing.StepBlock.recordTM, Turing.StepBlock.seekCfg, Turing.StepBlock.seekPart, Turing.StepBlock.seekRed, Turing.StepBlock.seekRouterTM, Turing.StepBlock.seekTM, Turing.StepBlock.summaryTape, Turing.inputMoveAction, Turing.phaseHostTM, Turing.tapeAction}
\difficulty{4}
Let $k$ be a natural number, let $x$ be a finite Boolean input, let $d$ be a
configuration of the record machine, let $rank<3^{k+1}$, and let $b$ be a Boolean.
Suppose $d$ is in the record machine's initial state and the symbol under the
summary-tape head is $b$. If $b$ is true, one step of the seek machine advances the
summary head by one and enters a request control carrying $rank$, the index $3^{k+1}-1$,
and the flag false. If $b$ is false, it leaves the configuration's tapes and head
positions unchanged and enters the done control when $rank=0$, or a request control
carrying $rank$, $rank-1$, and the flag true otherwise.
\end{lemma}

\begin{lemma}[Search past a ranked list of records]
\lean{Turing.StepBlock.seekTM_rank}
\label{Turing.StepBlock.seekTM_rank}
\leanok
\uses{Turing.Action, Turing.Cfg, Turing.Cfg.mapState, Turing.Cfg.workTapeSymbols, Turing.MultiTapeTM.runFrom, Turing.MultiTapeTM.runFrom_succ_eq_step, Turing.StepBlock.RecordState, Turing.StepBlock.recordTM, Turing.StepBlock.seekCfg, Turing.StepBlock.seekCost, Turing.StepBlock.seekTM, Turing.StepBlock.seekTM_many, Turing.StepBlock.seekTM_stateStep, Turing.StepBlock.summaryTape, Turing.StepBlock.tableTape, Turing.actionBitsK}
\difficulty{5}
Let $k,n$ be nonnegative integers, let $x$ be a finite binary input, let $A$ be a list
of transition actions, and let $d$ be a configuration. Let $r$ be a control rank such
that $|A|=r$, and suppose $d$ is in the recording machine's initial state, the
summary-tape head scans a blank, and the table tape stores the bit encoding of $A$
consecutively from its head. Starting the search machine in the state carrying rank $r$,
after $1+\operatorname{cost}(A)$ steps it is in the done control state, with its table
head advanced by the encoding length of $A$. It does not reach the done control state
earlier, and throughout this run its input-head position, tape contents, and output
remain unchanged, while every other work-tape head stays fixed.
\end{lemma}

\begin{definition}[Shifted work-tape head positions]
\lean{Turing.StepBlock.seekShift}
\label{Turing.StepBlock.seekShift}
\leanok
\uses{Turing.Cfg, Turing.StepBlock.RecordState, Turing.StepBlock.summaryTape, Turing.StepBlock.tableTape}
Let $k$ be a natural number, let $x$ be a finite Boolean input string, and let $d$ be a
configuration with $k+2$ work tapes, Boolean symbols, and record states. For natural
numbers $t$ and $s$, the shifted configuration preserves all data except the work-tape
head positions: the head on the table tape is moved by $t$, the head on the summary tape
is moved by $s$, and every other work-tape head remains in place.
\end{definition}

\begin{lemma}[Grouped table search]
\lean{Turing.StepBlock.seekTM_groups}
\label{Turing.StepBlock.seekTM_groups}
\leanok
\uses{Turing.Action, Turing.Cfg, Turing.Cfg.mapState, Turing.Cfg.workTapeSymbols, Turing.MultiTapeTM.runFrom, Turing.MultiTapeTM.runFrom_add, Turing.MultiTapeTM.runFrom_succ_eq_step, Turing.StepBlock.RecordState, Turing.StepBlock.reads_append, Turing.StepBlock.recordTM, Turing.StepBlock.seekCfg, Turing.StepBlock.seekCost, Turing.StepBlock.seekShift, Turing.StepBlock.seekTM, Turing.StepBlock.seekTM_many, Turing.StepBlock.seekTM_rank, Turing.StepBlock.seekTM_stateStep, Turing.StepBlock.summaryTape, Turing.StepBlock.summary_ne_table, Turing.StepBlock.tableTape, Turing.actionBitsK}
\difficulty{6}
Let $k,n$ be nonnegative integers, let $x$ be a finite binary input, and let $G$ be a
list of groups of transition actions, each group having length $3^{k+1}$. Let $R$ be a
list of remaining actions whose length is the rank $r$ of a control state, and put
$A=G\mathbin{\|}R$, where $G\mathbin{\|}R$ means the groups flattened and then
concatenated with $R$. Suppose a configuration $d$ is in the initial recording state,
its summary tape at offset $s$ from its head contains $1$ exactly when $s<|G|$ for every
$0\le s\le |G|$, and its table tape stores the bit encoding of $A$ consecutively from
its head. Starting the search machine in the control state carrying rank $r$, after
$T=|G|+1+\operatorname{cost}(A)$ steps it reaches the done control state, with the table
and summary heads shifted forward respectively by the length of the encoding of $A$ and
by $|G|$. It does not reach the done control state at any earlier time. Throughout times
$u\le T$, the input-head position, tape contents, and output remain those of $d$, and
every other work-tape head stays fixed.
\end{lemma}

\begin{lemma}[Initial observation in the search machine]
\lean{Turing.StepBlock.seekTM_observe}
\label{Turing.StepBlock.seekTM_observe}
\leanok
\uses{Turing.Cfg, Turing.Cfg.inputSymbol, Turing.Cfg.mapState, Turing.Cfg.workTapeSymbols, Turing.MultiTapeTM.runFrom, Turing.MultiTapeTM.step, Turing.StepBlock.RecordState, Turing.StepBlock.readRank, Turing.StepBlock.recordTM, Turing.StepBlock.seekCfg, Turing.StepBlock.seekShift, Turing.StepBlock.seekTM, Turing.StepBlock.summaryTape, Turing.StepBlock.summary_ne_table, Turing.StepBlock.tableTape, Turing.step_of_tapeAction_move}
\difficulty{5}
Let $k$ be a nonnegative integer, let $x$ be a finite binary input, and let $d$ be a
configuration whose state is the recording machine's initial state. Let $r$ be the rank
determined by the symbol under the input head and the symbols under the first $k$
work-tape heads of $d$. Starting the search machine in its observe control state, after
two steps it is in the state carrying rank $r$, with the summary-tape head advanced by
two cells. It does not reach the done state at either earlier time, and through time two
its input-head position, tape contents, and output are unchanged, while all work-tape
heads other than the summary-tape head stay fixed.
\end{lemma}

\begin{definition}[Actions from selected reads]
\lean{Turing.StepBlock.readActions}
\label{Turing.StepBlock.readActions}
\leanok
\uses{Turing.Action, Turing.workReadsK}
Let $k$ and $n$ be natural numbers, let $q$ be one of $n+1$ states, and let $tr$ be a
transition function for a $k$-tape machine with Boolean symbols and states in that set.
The list of read actions consists of the transitions at state $q$ for the prescribed
work-tape reads with a blank input symbol, followed by those for the same work-tape
reads with input symbol false, and then those with input symbol true.
\end{definition}

\begin{lemma}[Length of the transition-action list]
\lean{Turing.StepBlock.readActions_length}
\label{Turing.StepBlock.readActions_length}
\leanok
\uses{Turing.Action, Turing.StepBlock.length_workReadsK, Turing.StepBlock.readActions}
\difficulty{3}
Let $k,n\in\mathbb{N}$, let $\mathrm{tr}$ be a transition function assigning an action
to each state, input read, and work-tape read pattern, and let $q$ be a state in
$\{0,\ldots,n\}$. The list of transition actions at $q$, formed by listing the actions
for each of the three possible input reads and all work-tape read patterns, has length
$3^{k+1}$.
\end{lemma}

\begin{lemma}[Indexed transition actions]
\lean{Turing.StepBlock.readActions_at}
\label{Turing.StepBlock.readActions_at}
\leanok
\uses{Turing.Action, Turing.StepBlock.inputRank, Turing.StepBlock.length_workReadsK, Turing.StepBlock.mem_workReadsK, Turing.StepBlock.readActions, Turing.StepBlock.readActions_length, Turing.StepBlock.readRank, Turing.StepBlock.workRank, Turing.workReadsK}
\difficulty{5}
Let $k,n\in\mathbb{N}$, let $\mathrm{tr}$ assign an action to each state in
$\{0,\ldots,n\}$, input read in $\{\text{blank},0,1\}$, and work-tape read pattern, and
let $q$ be a state in $\{0,\ldots,n\}$. For any input read
$\mathrm{inp}\in\{\text{blank},0,1\}$ and any work-tape read pattern $w$, the action at
index $3^k\,\mathrm{rank}(\mathrm{inp})+\mathrm{rank}(w)$ in the list of all transition
actions at state $q$ is $\mathrm{tr}(q,\mathrm{inp},w)$, where the input ranks are
blank, $0$, $1$ in that order and the work-tape patterns are listed in the recursively
defined order.
\end{lemma}

\begin{lemma}[Grouping transition table entries by state]
\lean{Turing.StepBlock.tableBitsK_groups}
\label{Turing.StepBlock.tableBitsK_groups}
\leanok
\uses{Turing.Action, Turing.StepBlock.readActions, Turing.actionBitsK, Turing.tableBitsK}
\difficulty{1}
Let $k,n\in\mathbb{N}$, and let a transition table assign an action to each state in
$\{0,\ldots,n\}$, each input read in $\{\text{blank},0,1\}$, and each tuple of $k$
work-tape reads in $\{\text{blank},0,1\}^k$. The bit string encoding the table is the
concatenation, in state order, of the action encodings for each state, with that state's
entries ordered by input read and then by work-tape read tuple.
\end{lemma}

\begin{lemma}[Cutting a transition table into read blocks]
\lean{Turing.StepBlock.tableBitsK_cut}
\label{Turing.StepBlock.tableBitsK_cut}
\leanok
\uses{Turing.Action, Turing.StepBlock.readActions, Turing.StepBlock.readActions_at, Turing.StepBlock.readActions_length, Turing.StepBlock.readRank, Turing.StepBlock.tableBitsK_groups, Turing.actionBitsK, Turing.tableBitsK}
\difficulty{5}
Let $k,n\in\mathbb{N}$, and let a transition table assign an action to each state in
$\{0,\ldots,n\}$, each input read in $\{\text{blank},0,1\}$, and each tuple of $k$
work-tape reads in $\{\text{blank},0,1\}^k$. For any state $q$, input read $a$, and
work-tape read tuple $w$, the bit string encoding the full table can be decomposed into
the encodings of $q$ complete earlier state blocks, then $r$ read entries within state
$q$, then the encoding of the action at $(q,a,w)$, followed by a tail. Here there are
$q$ earlier blocks, each containing $3^{k+1}$ actions, and $r$ is the rank of $(a,w)$ in
the ordering that first orders input reads and then work-tape read tuples.
\end{lemma}

\begin{lemma}[Encoded Length of a List of Actions]
\lean{Turing.StepBlock.records_length}
\label{Turing.StepBlock.records_length}
\leanok
\uses{Turing.Action, Turing.actionBitsK, Turing.codeActionK_length}
\difficulty{4}
For every $k,n\in\mathbb N$ and every finite list of actions for a $k$-tape machine with
states indexed by $\{0,\ldots,n\}$, the number of actions is at most the length of the
concatenation of their binary encodings.
\end{lemma}

\begin{lemma}[Bound on work-tape head displacement]
\lean{Turing.MultiTapeTM.abs_workTapePos_runFrom_sub_le}
\label{Turing.MultiTapeTM.abs_workTapePos_runFrom_sub_le}
\leanok
\uses{Turing.Cfg, Turing.MultiTapeTM, Turing.MultiTapeTM.runFrom, Turing.MultiTapeTM.runFrom_succ_eq_step', Turing.MultiTapeTM.workTapePos_step_le}
\difficulty{4}
Let $M$ be a multi-tape Turing machine with $m$ work tapes and Boolean symbols, let $d$
be a configuration on input $x$, let $t\in\bbn$, and let $j<m$ be a work-tape index. The
absolute displacement of the head of tape $j$ after $t$ steps from its position in $d$
is at most $t$.
\end{lemma}

\begin{lemma}[Stationary input head]
\lean{Turing.MultiTapeTM.inputPos_runFrom_of_inputTape_eq_zero}
\label{Turing.MultiTapeTM.inputPos_runFrom_of_inputTape_eq_zero}
\leanok
\uses{Turing.Cfg, Turing.MultiTapeTM, Turing.MultiTapeTM.runFrom, Turing.MultiTapeTM.step}
\difficulty{4}
Let $M$ be a multi-tape Turing machine with Boolean tape symbols and state set $S$, and
let $x$ be a Boolean input word. Suppose every transition of $M$, for every state and
every scanned input and work-tape symbol, prescribes zero movement of the input head.
Then for every configuration $d$ on $x$ and every $t\in\mathbb{N}$, the input-head
position after running $M$ from $d$ for $t$ steps equals its initial position in $d$.
\end{lemma}

\begin{lemma}[Length of flattened equal-length groups]
\lean{Turing.StepBlock.groups_length}
\label{Turing.StepBlock.groups_length}
\leanok
\uses{Turing.Action}
\difficulty{4}
Let $k,n$ be natural numbers and let $G$ be a finite list of lists of actions for $k$
Boolean work tapes and successor states in $\mathrm{Fin}(n+1)$. If every list in $G$ has
length $3^{k+1}$, then the concatenation of all lists in $G$ has length $|G|\,3^{k+1}$.
\end{lemma}

\begin{lemma}[Lookup from the initial search state]
\lean{Turing.StepBlock.seekTM_lookup}
\label{Turing.StepBlock.seekTM_lookup}
\leanok
\uses{Turing.Action, Turing.Cfg, Turing.Cfg.inputSymbol, Turing.Cfg.mapState, Turing.Cfg.workTapeSymbols, Turing.MultiTapeTM.runFrom, Turing.MultiTapeTM.runFrom_add, Turing.StepBlock.RecordState, Turing.StepBlock.groups_length, Turing.StepBlock.readRank, Turing.StepBlock.recordTM, Turing.StepBlock.records_length, Turing.StepBlock.seekCfg, Turing.StepBlock.seekCost, Turing.StepBlock.seekShift, Turing.StepBlock.seekTM, Turing.StepBlock.seekTM_groups, Turing.StepBlock.seekTM_observe, Turing.StepBlock.summaryTape, Turing.StepBlock.summary_ne_table, Turing.StepBlock.tableTape, Turing.actionBitsK}
\difficulty{5}
Let $k,n$ be nonnegative integers, let $x$ be a finite binary input, let $G$ be a list
of groups of transition actions, and let $R$ be a list of remaining actions. Suppose
every group has length $3^{k+1}$, the length of $R$ equals the rank read from the input
symbol and the symbols under the first $k$ work-tape heads of a configuration $d$, and
$d$ is in the recording machine's initial state. Suppose the summary tape, starting two
cells beyond its head, marks precisely the group positions: at offset $2+s$ it contains
$1$ exactly when $s<|G|$, for $0\le s\le |G|$. Suppose the table tape stores the
encoding of the flattened groups followed by $R$, consecutively from its head. Put
$A=G\mathbin{\|}R$ and $T=|G|+3+\operatorname{cost}(A)$. Starting the search machine in
its observe control state, after $T$ steps it reaches the done control state, with its
table and summary heads shifted forward respectively by the encoding length of $A$ and
by $|G|+2$. It does not reach the done control state earlier, and through time $T$ its
input-head position, tape contents, and output remain unchanged, while every other
work-tape head stays fixed. Moreover, $T\le 5$ times the encoding length of $A$, plus
$3$.
\end{lemma}

\begin{lemma}[Reading a unary word]
\lean{Turing.StepBlock.bufferTape_unary}
\label{Turing.StepBlock.bufferTape_unary}
\leanok
\uses{Turing.FinTM.bufferTape, Turing.FinTM.bufferTape_nat}
\difficulty{4}
Let $q,r$ be natural numbers with $r\le q$. In the buffer tape storing $q$ copies of $1$
followed by $0$, cell $r$ contains $1$ if $r<q$ and contains $0$ if $r=q$.
\end{lemma}

\begin{definition}[Summary tape cell contents]
\lean{Turing.StepBlock.summaryCells}
\label{Turing.StepBlock.summaryCells}
\leanok
\uses{Turing.FinTM.bufferTape}
For an optional Boolean summary symbol, a finite Boolean word, and an integer position
$z$, define the cell contents to be the summary symbol when $z=0$, and otherwise to be
the buffer-tape contents of the word at position $z-2$.
\end{definition}

\begin{lemma}[Shifted summary-cell values]
\lean{Turing.StepBlock.summaryCells_shift}
\label{Turing.StepBlock.summaryCells_shift}
\leanok
\uses{Turing.FinTM.bufferTape, Turing.StepBlock.summaryCells}
\difficulty{3}
Let $s\in\mathrm{Option}(\mathrm{Bool})$, let $w$ be a Boolean word, and let
$z\in\mathbb Z$ satisfy $-1\le z$. The summary-cell content at position $2+z$ equals the
buffer content of $w$ at position $z$.
\end{lemma}

\begin{lemma}[Updating the summary cell]
\lean{Turing.StepBlock.summaryCells_update}
\label{Turing.StepBlock.summaryCells_update}
\leanok
\uses{Turing.StepBlock.summaryCells}
\difficulty{4}
Let $s,n\in\mathrm{Option}(\mathrm{Bool})$ and let $w$ be a Boolean word. Replacing the
summary-cell function for $s$ and $w$ at position zero by $n$ gives the summary-cell
function for $n$ and $w$.
\end{lemma}

\begin{lemma}[Blank cells outside the stored word]
\lean{Turing.StepBlock.bufferTape_outside}
\label{Turing.StepBlock.bufferTape_outside}
\leanok
\uses{Turing.FinTM.bufferTape}
\difficulty{4}
Let $w$ be a finite Boolean word and $z$ an integer. If $z$ is not a nonnegative integer
strictly less than the length of $w$, then the buffer tape storing $w$ is blank at cell
$z$.
\end{lemma}

\begin{lemma}[Copying the summary word region]
\lean{Turing.StepBlock.summaryCells_copy}
\label{Turing.StepBlock.summaryCells_copy}
\leanok
\uses{Turing.FinTM.bufferTape, Turing.StepBlock.bufferTape_outside, Turing.StepBlock.summaryCells}
\difficulty{5}
Let $s\in\mathrm{Option}(\mathrm{Bool})$ and let $w$ be a Boolean word. At each integer
cell $z$, assign the buffer content of $w$ shifted by two when $2\le z<2+|w|$, and
otherwise use the summary-cell content for $s$ and the empty word. The resulting
function equals the summary-cell function for $s$ and $w$.
\end{lemma}

\begin{lemma}[Clearing the summary word region]
\lean{Turing.StepBlock.summaryCells_clear}
\label{Turing.StepBlock.summaryCells_clear}
\leanok
\uses{Turing.StepBlock.summaryCells, Turing.StepBlock.summaryCells_copy}
\difficulty{4}
Let $s\in\{\mathrm{none},\mathrm{some\ false},\mathrm{some\ true}\}$ and let $w$ be a
Boolean word. At every integer cell $z$, the function that replaces the summary-cell
content by blank for $2\le z<2+|w|$ and otherwise retains the content associated to $s$
and the empty word agrees with the summary-cell function associated to $s$ and the empty
word.
\end{lemma}

\begin{lemma}[Selection of a transition-table entry]
\lean{Turing.StepBlock.seekTM_select}
\label{Turing.StepBlock.seekTM_select}
\leanok
\uses{Turing.Action, Turing.Cfg, Turing.Cfg.inputSymbol, Turing.Cfg.workTapeSymbols, Turing.FinTM.bufferTape, Turing.FinTM.bufferTape_nat, Turing.MultiTapeTM.runFrom, Turing.StepBlock.RecordState, Turing.StepBlock.bufferTape_unary, Turing.StepBlock.groups_length, Turing.StepBlock.reads_append, Turing.StepBlock.recordTM, Turing.StepBlock.records_length, Turing.StepBlock.seekCfg, Turing.StepBlock.seekCost, Turing.StepBlock.seekShift, Turing.StepBlock.seekTM, Turing.StepBlock.seekTM_lookup, Turing.StepBlock.summaryCells, Turing.StepBlock.summaryCells_shift, Turing.StepBlock.summaryTape, Turing.StepBlock.tableBitsK_cut, Turing.StepBlock.tableTape, Turing.actionBitsK, Turing.tableBitsK, Turing.unaryFin}
\difficulty{6}
Let $k,n$ be natural numbers, let $x$ be a finite Boolean input, let $tr$ be a
transition function for $n+1$ states and $k$ work tapes, and let $q$ be one of those
states. Let $summary$ be an optional bit and $d$ a configuration of the record machine.
Suppose $d$ is in the record machine's initial state, its table tape stores the
transition table encoding of $tr$ starting at cell zero, its table head is at zero, its
summary tape stores $summary$ at cell zero followed by the unary encoding of $q$, and
its summary head is at zero. For the action selected by $tr$ on state $q$, the symbol
currently under the input head, and the symbols under the work-tape heads, there are bit
strings $pre,tail$ and a positive integer $T$ such that the table encoding is $pre$
followed by that action's encoding and then $tail$, the unary state index is no larger
than the length of $pre$, and $T$ is at most five times the table-encoding length plus
three. After $T$ steps from the seek machine's observe configuration, the machine is in
its done control state, with the table head advanced by $|pre|$ and the summary head
advanced by $q+2$. Before step $T$ it is never in the terminal control state; through
step $T$, the input-head position, tape contents, output, and all work-tape head
positions other than the table and summary heads remain unchanged.
\end{lemma}

\begin{definition}[Two-move machine]
\lean{Turing.StepBlock.twoMoveTM}
\label{Turing.StepBlock.twoMoveTM}
\leanok
\uses{Turing.MultiTapeTM, Turing.inputMoveAction, Turing.tapeAction}
Let $m$ be a natural number, let $i$ be one of its tape indices, and let the requested
movement be left, stay, or right. The machine has three states and starts in state $0$.
In either of states $0$ or $1$, it performs the specified action on tape $i$ and moves
to the next state; in state $2$, it performs the input-movement action for that state.
\end{definition}

\begin{lemma}[Running the two-move machine]
\lean{Turing.StepBlock.twoMoveTM_run}
\label{Turing.StepBlock.twoMoveTM_run}
\leanok
\uses{Turing.Cfg, Turing.MultiTapeTM.runFrom, Turing.MultiTapeTM.runFrom_succ_eq_step', Turing.StepBlock.twoMoveTM, Turing.step_of_tapeAction_move}
\difficulty{4}
Let $m\in\mathbb{N}$, let $i$ be one of $m$ work tapes, and choose a head movement
$\delta\in\{-1,0,1\}$. For any input string $x$ and configuration $d$ whose state is
$0$, after every $u\in\mathbb{N}$ with $u\le 2$, the machine that in states $0$ and $1$
moves tape $i$ by $\delta$ and advances to the next state, and otherwise makes no change
to the tapes and stays in its state, reaches the configuration obtained from $d$ by
setting the state to $u$ and the head position of tape $i$ to its original position plus
$u\delta$; all other configuration components are unchanged.
\end{lemma}

\begin{definition}[Finite set of control keys]
\lean{Turing.StepBlock.Key}
\label{Turing.StepBlock.Key}
\leanok
The type of keys has the following constructors: entry, lookup, state tail, state home,
clear state, to summary, record, to state, successor, halt; table tail and table home,
each indexed by a Boolean live flag; state gap and state restore, each indexed by a
Boolean live flag; and finish, indexed by a Boolean live flag.
\end{definition}

\begin{definition}[State types for machine phases]
\lean{Turing.StepBlock.Ph}
\label{Turing.StepBlock.Ph}
\leanok
\uses{Turing.SeekPhase, Turing.StepBlock.Key, Turing.StepBlock.RecordState, Turing.StepBlock.SeekPh, Turing.StepBlock.SuccPhase, Turing.SweepPhase}
For each natural number $k$, assign a state type to every control key as follows: the
entry and halt keys have a singleton state type; the lookup key has states consisting of
a seek key together with the corresponding seek-phase state; the state-tail, state-home,
table-tail, table-home, state-gap, and state-restore keys have seek-phase states; the
clear-state key has sweep-phase states; the to-summary, to-state, and finish keys have
states in a three-element set; the record key has record states of width $4k+4$; and the
successor key has successor-phase states.
\end{definition}

\begin{definition}[Machines for each control phase]
\lean{Turing.StepBlock.part}
\label{Turing.StepBlock.part}
\leanok
\uses{Turing.MultiTapeTM, Turing.StepBlock.Key, Turing.StepBlock.Ph, Turing.StepBlock.recordTM, Turing.StepBlock.seekTM, Turing.StepBlock.succFieldTM, Turing.StepBlock.summaryTape, Turing.StepBlock.tableTape, Turing.StepBlock.twoMoveTM, Turing.clearTM, Turing.inputMoveAction, Turing.skipBackTM, Turing.skipTM}
For each natural number $k$ and each control key, assign a deterministic Turing machine
over the Boolean alphabet with $k+2$ work tapes and the state type assigned to that key.
The entry and halt machines leave all tapes stationary; lookup uses the seek machine;
state-tail and state-gap use the forward skip machine on the summary tape, while
state-home and state-restore use its return-pass version; clear-state clears the summary
tape; to-summary and finish make two leftward moves on that tape, and to-state makes two
rightward moves. The record key uses the record machine, the successor key uses the
successor-field machine copying to the summary tape, and table-tail and table-home use
the forward and return-pass skip machines on the table tape, respectively.
\end{definition}

\begin{definition}[Live or halted exit state]
\lean{Turing.StepBlock.exit}
\label{Turing.StepBlock.exit}
\leanok
\uses{Turing.StepBlock.Ph}
For each natural number $k$ and Boolean flag $b$, define the exit state to be the entry
key with its unique state when $b$ is true, and the halt key with its unique state when
$b$ is false.
\end{definition}

\begin{definition}[Control-phase reduction rule]
\lean{Turing.StepBlock.red}
\label{Turing.StepBlock.red}
\leanok
\uses{Turing.StepBlock.Key, Turing.StepBlock.Ph, Turing.StepBlock.exit, Turing.StepBlock.recordTM, Turing.StepBlock.seekTM}
For each natural number $k$, define a partial reduction from a control key and its phase
state to a new control key and phase state as follows. Entry advances to lookup; lookup
advances to state-tail precisely when its state is the empty-key done state; state-tail
and state-home advance, on their done state, to state-home and clear-state,
respectively; clear-state advances from done to to-summary at position $0$; to-summary
advances from position $2$ to record; record advances from its halting state to to-state
at position $0$; to-state advances from position $2$ to successor at the tag phase; a
completed successor field advances to table-tail, which on completion advances to
table-home, then to state-gap, then to state-restore, and finally to finish, preserving
the live flag at each step; finish at position $2$ advances to the exit state determined
by that flag. In every other case the reduction is undefined.
\end{definition}

\begin{definition}[One-round composite machine]
\lean{Turing.StepBlock.roundTM}
\label{Turing.StepBlock.roundTM}
\leanok
\uses{Turing.MultiTapeTM, Turing.StepBlock.Ph, Turing.StepBlock.part, Turing.StepBlock.red, Turing.phaseHostTM}
For each natural number $k$, define a deterministic Turing machine over the Boolean
alphabet with $k+2$ work tapes and states consisting of a control key together with a
state of that key’s assigned phase machine. Its initial state is the entry key with its
unique state, and its transitions combine the phase machines with the control-phase
reduction rule.
\end{definition}

\begin{definition}[Finite one-round machine]
\lean{Turing.StepBlock.roundFinTM}
\label{Turing.StepBlock.roundFinTM}
\leanok
\uses{Turing.FinTM, Turing.StepBlock.Ph, Turing.StepBlock.roundTM}
For each natural number $k$, define a finite Turing machine over the Boolean alphabet
with $k+2$ work tapes, state set consisting of a control key together with a state of
that key’s assigned phase machine, and underlying machine the one-round composite
machine.
\end{definition}

\begin{definition}[Augmented record configuration]
\lean{Turing.StepBlock.roundData}
\label{Turing.StepBlock.roundData}
\leanok
\uses{Turing.Cfg, Turing.FinTM.bufferTape, Turing.StepBlock.RecordState, Turing.StepBlock.recordTM, Turing.StepBlock.summaryCells}
Let $k,n$ be natural numbers, let $x$ be a Boolean input list, and let $c$ be a
configuration with $k$ work tapes, Boolean alphabet, and states in the finite set of
size $n+1$. Given a Boolean table word, an optional Boolean summary, a Boolean word,
integer table and state head positions, and an output prefix, define a configuration
with $k+2$ work tapes and record states by setting its state to the initial
record-machine state, retaining the input-head position of $c$, using the first $k$ work
tapes and head positions of $c$, placing the table buffer on the next tape and the
summary cells on the final tape with the specified head positions, and setting the
output to the given prefix.
\end{definition}

\begin{lemma}[Distinctness of an Original and Summary Tape]
\lean{Turing.StepBlock.mirror_ne_summary}
\label{Turing.StepBlock.mirror_ne_summary}
\leanok
\uses{Turing.StepBlock.summaryTape}
\difficulty{3}
For every natural number $k$ and every $i\in\{0,\ldots,k-1\}$, the tape indexed by $i$
among the original $k$ tapes is distinct from the additional summary tape.
\end{lemma}

\begin{lemma}[Shift of Round-Configuration Heads]
\lean{Turing.StepBlock.roundData_shift}
\label{Turing.StepBlock.roundData_shift}
\leanok
\uses{Turing.Cfg, Turing.StepBlock.mirror_ne_summary, Turing.StepBlock.mirror_ne_table, Turing.StepBlock.roundData, Turing.StepBlock.seekShift, Turing.StepBlock.summaryTape, Turing.StepBlock.tableTape}
\difficulty{4}
Let $k,n\in\mathbb N$, let $x$ be a finite binary input, let $c$ be a configuration of a
machine with $k$ work tapes and states indexed by $\{0,\ldots,n\}$, and let a table
word, a summary symbol, a state word, two integer head offsets, and an output list be
given. Shifting the table and state head positions of the round configuration by natural
numbers $a$ and $b$ yields exactly the round configuration whose table and state head
positions have been increased by $a$ and $b$, respectively.
\end{lemma}

\begin{definition}[One-round head-position window]
\lean{Turing.StepBlock.RoundWindow}
\label{Turing.StepBlock.RoundWindow}
\leanok
\uses{Turing.Action, Turing.Action.apply, Turing.Cfg}
Let $c$ be a configuration with $k$ work tapes, Boolean alphabet, and states in a set of
type $S$, on input $x$, and let $a$ be an action for a machine with $k$ work tapes on
that input. A configuration $d$ with $k+2$ work tapes on the same input satisfies the
one-round window condition if its input-head position is either the position in $c$ or
the position after applying $a$ to $c$, and, for every one of the original $k$ work
tapes, its corresponding head position is either the position in $c$ or the position
after applying $a$ to $c$.
\end{definition}

\begin{lemma}[Round Data Window]
\lean{Turing.StepBlock.roundData_window}
\label{Turing.StepBlock.roundData_window}
\leanok
\uses{Turing.Action, Turing.Action.apply, Turing.Cfg, Turing.StepBlock.RoundWindow, Turing.StepBlock.roundData}
\difficulty{3}
Let $c$ be a configuration of a machine with $k$ work tapes, let $a$ be a one-step
action, and let $d$ be the configuration with two additional work tapes constructed from
either $c$ or the result of applying $a$ to $c$, according to a Boolean choice. The
input-head position of $d$ is either that of $c$ or that after applying $a$, and each of
the original $k$ work-head positions in $d$ is either its position in $c$ or its
position after applying $a$.
\end{lemma}

\begin{definition}[Round execution span]
\lean{Turing.StepBlock.RoundSpan}
\label{Turing.StepBlock.RoundSpan}
\leanok
\uses{Turing.Action, Turing.Cfg, Turing.MultiTapeTM.runFrom, Turing.StepBlock.Ph, Turing.StepBlock.RoundWindow, Turing.StepBlock.roundTM}
For a $k$-tape configuration $c$ over Boolean work symbols and states in
$	ext{Fin}(n+1)$, an action $a$ of the same machine, two configurations $d,e$ of the
compiled machine with $k+2$ work tapes, and a time $T\in\mathbb{N}$, the round-span
condition holds when the compiled round machine started at $d$ reaches $e$ after exactly
$T$ steps, and at every time $u<T$ its state is neither the entry state nor the halt
state and its configuration satisfies the round-window condition determined by $c$ and
$a$.
\end{definition}

\begin{lemma}[Concatenating round spans]
\lean{Turing.StepBlock.RoundSpan.add}
\label{Turing.StepBlock.RoundSpan.add}
\leanok
\uses{Turing.Action, Turing.Cfg, Turing.MultiTapeTM.runFrom_add, Turing.StepBlock.Ph, Turing.StepBlock.RoundSpan}
\difficulty{4}
Let $k,n$ be natural numbers, let $x$ be a finite Boolean word, let $c$ be a
configuration of a $k$-tape machine with states in $\mathrm{Fin}(n+1)$, and let $a$ be
an action for that machine. Suppose configurations $d,e,f$ of the associated
$(k+2)$-tape round machine satisfy the following: the run from $d$ reaches $e$ after $T$
steps, and at every earlier time its state is present and is neither the entry state nor
the halt state, while its input-head position and each of its first $k$ work-tape head
positions agree with either the corresponding position in $c$ or the position after
applying $a$ to $c$; and the run from $e$ reaches $f$ after $U$ steps with those same
intermediate-time conditions. Then the run from $d$ reaches $f$ after $T+U$ steps and
satisfies those conditions at every time strictly before $T+U$.
\end{lemma}

\begin{lemma}[A redirected phase forms a round span]
\lean{Turing.StepBlock.roundTM_phase}
\label{Turing.StepBlock.roundTM_phase}
\leanok
\uses{Turing.Action, Turing.Cfg, Turing.Cfg.mapState, Turing.MultiTapeTM.runFrom, Turing.MultiTapeTM.runFrom_add, Turing.MultiTapeTM.runFrom_of_halt, Turing.StepBlock.Key, Turing.StepBlock.Ph, Turing.StepBlock.RoundSpan, Turing.StepBlock.RoundWindow, Turing.StepBlock.part, Turing.StepBlock.red, Turing.phaseHostTM_run}
\difficulty{5}
Let $k,n\in\mathbb{N}$, let $x$ be a finite binary input, let $c$ be a configuration of
a $k$-tape machine with states indexed by $\{0,\ldots,n\}$, and let $a$ be one of its
actions. Choose a phase label $j$ other than the entry or halt label, configurations
$d,e$ in that phase, and a time $T$. Suppose the phase machine runs from $d$ to $e$ in
$T$ steps; the state of $e$ is $s$; the phase redirection at $(j,s)$ produces a next
labelled phase; no state encountered at any time $u<T$ redirects; and the round-window
property holds at every phase-machine time $u\le T$. Then the full round machine runs
from $d$ embedded under label $j$ to $e$ embedded in the redirected next phase in $T+1$
steps, with the round-span property.
\end{lemma}

\begin{definition}[Configuration at a designated state]
\lean{Turing.StepBlock.atKey}
\label{Turing.StepBlock.atKey}
\leanok
\uses{Turing.Cfg, Turing.Cfg.mapState, Turing.StepBlock.Key, Turing.StepBlock.Ph, Turing.StepBlock.RecordState}
Given a configuration $d$ of a machine with $k+2$ work tapes, whose states are record
states, a key $q$ and a phase value $p$ of the type assigned to that key determine a
configuration with the same input-head position, work-tape contents and head positions,
and output as $d$, but with its state replaced by the pair $(q,p)$.
\end{definition}

\begin{lemma}[Preservation of the round window]
\lean{Turing.StepBlock.roundWindow_frame}
\label{Turing.StepBlock.roundWindow_frame}
\leanok
\uses{Turing.Action, Turing.Cfg, Turing.StepBlock.RoundWindow}
\difficulty{3}
Let $c$ be a configuration with $k$ work tapes, let $a$ be an action on $c$, and let $d$
be a configuration with two additional work tapes. Suppose the input-head position of
$d$ is either its position in $c$ or its position after applying $a$, and each of its
first $k$ work-tape heads is likewise at the corresponding position in $c$ or after
applying $a$. If a configuration $e$ has the same input-head position as $d$ and the
same positions for those first $k$ work-tape heads, then $e$ satisfies the same two
alternatives for each of these positions.
\end{lemma}

\begin{lemma}[State-Word Cells in the Summary Tape]
\lean{Turing.StepBlock.roundData_stateCell}
\label{Turing.StepBlock.roundData_stateCell}
\leanok
\uses{Turing.Cfg, Turing.FinTM.bufferTape, Turing.StepBlock.roundData, Turing.StepBlock.summaryCells_shift, Turing.StepBlock.summaryTape}
\difficulty{1}
Let $k,n\in\mathbb N$, let $x$ be a finite binary input, let $c$ be a configuration of a
machine with $k$ work tapes and states indexed by $\{0,\ldots,n\}$, and let a table
word, summary symbol, state word, integer head positions, and output list be given. For
every integer $z\ge -1$, the symbol on the summary tape of the resulting round
configuration at cell $2+z$ is the symbol stored by the buffer tape for the state word
at cell $z$.
\end{lemma}

\begin{lemma}[Preparing the record phase]
\lean{Turing.StepBlock.roundTM_prepare}
\label{Turing.StepBlock.roundTM_prepare}
\leanok
\uses{Turing.Action, Turing.Cfg, Turing.Cfg.mapState, Turing.FinTM.bufferTape, Turing.FinTM.bufferTape_nat, Turing.MultiTapeTM.inputPos_runFrom_of_inputTape_eq_zero, Turing.MultiTapeTM.runFrom, Turing.StepBlock.RoundSpan, Turing.StepBlock.RoundSpan.add, Turing.StepBlock.atKey, Turing.StepBlock.mirror_ne_summary, Turing.StepBlock.recordTM, Turing.StepBlock.red, Turing.StepBlock.roundData, Turing.StepBlock.roundData_stateCell, Turing.StepBlock.roundData_window, Turing.StepBlock.roundTM_phase, Turing.StepBlock.roundWindow_frame, Turing.StepBlock.summaryCells, Turing.StepBlock.summaryCells_clear, Turing.StepBlock.summaryTape, Turing.StepBlock.twoMoveTM, Turing.StepBlock.twoMoveTM_run, Turing.clearTM, Turing.clearTM_run_ofCfg, Turing.inputMoveAction, Turing.skipBackTM, Turing.skipBackTM_run_ofCfg, Turing.skipTM, Turing.skipTM_run_ofCfg, Turing.tapeAction, Turing.unaryFin}
\difficulty{6}
Let $k,n\in\mathbb{N}$, let $x$ be a finite binary input, let $c$ be a configuration of
a $k$-tape machine with states indexed by $\{0,\ldots,n\}$, let $a$ be an action, and
choose a state index $q$, a binary table, an optional summary bit, an integer position
$p$, and an output prefix. The round machine runs from a state-tail configuration whose
summary tape stores the unary encoding of $q$ at position $q+2$ and whose table head is
at $p$, to the record-phase initial configuration with the same table and summary but
with the summary tape empty and its head at $0$. The run takes $3q+16$ steps and
satisfies the round-span property.
\end{lemma}

\begin{lemma}[Reading a word stored between two lists]
\lean{Turing.StepBlock.bufferTape_reads}
\label{Turing.StepBlock.bufferTape_reads}
\leanok
\uses{Turing.FinTM.bufferTape, Turing.FinTM.bufferTape_nat, Turing.StepBlock.reads_append}
\difficulty{4}
Let $b,w,a$ be finite Boolean words. For every natural number $r<|w|$, the buffer tape
storing their concatenation $bwa$ contains the $r$-th bit of $w$ at cell $|b|+r$.
\end{lemma}

\begin{lemma}[Processing one action payload]
\lean{Turing.StepBlock.roundTM_payload}
\label{Turing.StepBlock.roundTM_payload}
\leanok
\uses{Turing.Action, Turing.Action.apply, Turing.Cfg, Turing.Cfg.mapState, Turing.Cfg.workTapeSymbols, Turing.MultiTapeTM.runFrom, Turing.StepBlock.RecordState, Turing.StepBlock.RoundSpan, Turing.StepBlock.RoundSpan.add, Turing.StepBlock.actionBitsK_reads, Turing.StepBlock.afterTag, Turing.StepBlock.applyPayload, Turing.StepBlock.atKey, Turing.StepBlock.bufferTape_reads, Turing.StepBlock.mirror_ne_summary, Turing.StepBlock.mirror_ne_table, Turing.StepBlock.outputSummary, Turing.StepBlock.outputSummary_append, Turing.StepBlock.readCfg, Turing.StepBlock.recordTM, Turing.StepBlock.recordTM_effect, Turing.StepBlock.recordTM_execute, Turing.StepBlock.recordWidth, Turing.StepBlock.red, Turing.StepBlock.roundData, Turing.StepBlock.roundTM_phase, Turing.StepBlock.succFieldCfg, Turing.StepBlock.succFieldTM, Turing.StepBlock.succFieldTM_cut, Turing.StepBlock.succFieldTM_heads, Turing.StepBlock.succFieldTM_run, Turing.StepBlock.summaryCells, Turing.StepBlock.summaryCells_copy, Turing.StepBlock.summaryCells_update, Turing.StepBlock.summaryEmit, Turing.StepBlock.summaryTape, Turing.StepBlock.summary_ne_table, Turing.StepBlock.tableTape, Turing.StepBlock.twoMoveTM, Turing.StepBlock.twoMoveTM_run, Turing.actionBitsK, Turing.optStateBits, Turing.unaryFin}
\difficulty{6}
Let $k,n\in\mathbb{N}$, let $x$ be a finite binary input, let $c$ be a configuration of
a $k$-tape machine with states indexed by $\{0,\ldots,n\}$, and let $a$ be an action for
that machine. For any binary words
$\mathrm{table},\mathrm{before},\mathrm{after},\mathrm{pre}$ with
$\mathrm{table}=\mathrm{before}$ followed by the serialization of $a$ followed by
$\mathrm{after}$, the round machine runs from the record configuration with this table,
the summary of $c$'s output, and the heads at positions $|\mathrm{before}|$ and $0$, to
the table-tail configuration after the action. The state tape then contains the unary
encoding of the successor state (or is empty if there is none), the table head is at
$|\mathrm{before}|+|a|_{\mathrm{code}}$, and the run takes $|a|_{\mathrm{code}}+6$
steps, satisfying the round-window property throughout the run as specified.
\end{lemma}

\begin{lemma}[Buffer tape under a nonnegative shift]
\lean{Turing.StepBlock.bufferTape_suffix}
\label{Turing.StepBlock.bufferTape_suffix}
\leanok
\uses{Turing.FinTM.bufferTape, Turing.FinTM.bufferTape_nat}
\difficulty{3}
Let $b,w$ be finite Boolean words and let $z$ be a nonnegative integer. The buffer tape
storing $bw$, read at cell $|b|+z$, has the same content as the buffer tape storing $w$
at cell $z$.
\end{lemma}

\begin{lemma}[Restoring tapes at the end of a round]
\lean{Turing.StepBlock.roundTM_restore}
\label{Turing.StepBlock.roundTM_restore}
\leanok
\uses{Turing.Action, Turing.Action.apply, Turing.Cfg, Turing.Cfg.mapState, Turing.FinTM.bufferTape_nat, Turing.MultiTapeTM.inputPos_runFrom_of_inputTape_eq_zero, Turing.MultiTapeTM.runFrom, Turing.StepBlock.RoundSpan, Turing.StepBlock.RoundSpan.add, Turing.StepBlock.atKey, Turing.StepBlock.bufferTape_suffix, Turing.StepBlock.exit, Turing.StepBlock.mirror_ne_summary, Turing.StepBlock.mirror_ne_table, Turing.StepBlock.outputSummary, Turing.StepBlock.red, Turing.StepBlock.roundData, Turing.StepBlock.roundData_stateCell, Turing.StepBlock.roundData_window, Turing.StepBlock.roundTM_phase, Turing.StepBlock.roundWindow_frame, Turing.StepBlock.summaryTape, Turing.StepBlock.tableTape, Turing.StepBlock.twoMoveTM, Turing.StepBlock.twoMoveTM_run, Turing.inputMoveAction, Turing.skipBackTM, Turing.skipBackTM_run_ofCfg, Turing.skipTM, Turing.skipTM_run_ofCfg, Turing.tapeAction}
\difficulty{7}
Let $k,n\in\mathbb{N}$, let $x$ be a finite binary input, let $c$ be a configuration of
a $k$-tape machine with states indexed by $\{0,\ldots,n\}$, and let $a$ be an action.
Choose binary words
$\mathrm{table},\mathrm{before},\mathrm{tail},\mathrm{word},\mathrm{pre}$ and a Boolean
flag $\ell$, with $\mathrm{table}=\mathrm{before}$ followed by $\mathrm{tail}$. The
round machine runs from the table-tail configuration based on the post-action
configuration $a(c)$, with the table head at $|\mathrm{before}|$ and the summary head at
$2+|\mathrm{word}|$, to the configuration based on $a(c)$ with both auxiliary heads at
$0$ and the exit label determined by $\ell$. It takes
$|\mathrm{tail}|+|\mathrm{table}|+|\mathrm{word}|+15$ steps and satisfies the round-span
property.
\end{lemma}

\begin{lemma}[Lookup within one machine round]
\lean{Turing.StepBlock.roundTM_lookup}
\label{Turing.StepBlock.roundTM_lookup}
\leanok
\uses{Turing.Action, Turing.Cfg, Turing.Cfg.inputSymbol, Turing.Cfg.mapState, Turing.Cfg.workTapeSymbols, Turing.StepBlock.RoundSpan, Turing.StepBlock.atKey, Turing.StepBlock.mirror_ne_summary, Turing.StepBlock.mirror_ne_table, Turing.StepBlock.outputSummary, Turing.StepBlock.red, Turing.StepBlock.roundData, Turing.StepBlock.roundData_shift, Turing.StepBlock.roundData_window, Turing.StepBlock.roundTM_phase, Turing.StepBlock.roundWindow_frame, Turing.StepBlock.seekCfg, Turing.StepBlock.seekShift, Turing.StepBlock.seekTM, Turing.StepBlock.seekTM_select, Turing.StepBlock.summaryTape, Turing.StepBlock.tableTape, Turing.actionBitsK, Turing.tableBitsK, Turing.unaryFin}
\difficulty{6}
Let $k,n\in\mathbb{N}$, let $x$ be a finite binary input, let $c$ be a configuration of
a $k$-tape machine with states indexed by $\{0,\ldots,n\}$, let $q$ be such a state, and
let $\mathrm{tr}$ be a transition table assigning an action to each state, input symbol
(or blank), and tuple of scanned work-tape symbols. Let $a$ be the action prescribed by
this table at $q$ and the symbols scanned in $c$. There are binary words $b,t$ and a
time $T$ such that the serialization of the entire transition table is $b$ followed by
the serialization of $a$ followed by $t$, $q\le |b|$, and $T\le 5$ times the length of
the serialized table plus $4$; moreover, the round machine runs from the lookup
configuration with the table stored on its table tape, the output summary and unary
encoding of $q$ on its summary tape, both tape heads at $0$, and output prefix
unchanged, to the state-tail configuration after $T$ steps, with the stated round-window
property at every earlier step.
\end{lemma}

\begin{lemma}[Canonical Round Simulation]
\lean{Turing.StepBlock.roundTM_canonical}
\label{Turing.StepBlock.roundTM_canonical}
\leanok
\uses{Turing.Action, Turing.Action.apply, Turing.Cfg, Turing.Cfg.inputSymbol, Turing.Cfg.mapState, Turing.Cfg.workTapeSymbols, Turing.MultiTapeTM.runFrom, Turing.MultiTapeTM.runFrom_succ_eq_step, Turing.MultiTapeTM.step, Turing.StepBlock.RoundSpan.add, Turing.StepBlock.RoundWindow, Turing.StepBlock.actionBitsK_split, Turing.StepBlock.atKey, Turing.StepBlock.exit, Turing.StepBlock.outputSummary, Turing.StepBlock.roundData, Turing.StepBlock.roundData_window, Turing.StepBlock.roundTM, Turing.StepBlock.roundTM_lookup, Turing.StepBlock.roundTM_payload, Turing.StepBlock.roundTM_prepare, Turing.StepBlock.roundTM_restore, Turing.StepBlock.seekTM, Turing.actionBitsK, Turing.moveInputPos_zero, Turing.optStateBits, Turing.step_of_inputMoveAction, Turing.tableBitsK, Turing.unaryFin}
\difficulty{6}
Let $k,n$ be natural numbers, let $x$ be a Boolean input word, let $c$ be a
configuration with $k$ Boolean work tapes and $n+1$ possible states, and let $q$ be one
of those states. Let $\mathrm{tr}$ be a transition rule assigning an action to each
state, optional input bit, and tuple of optional work-tape bits; write $a$ for the
action selected by $\mathrm{tr}$ from $q$ and the symbols scanned in $c$. For any output
prefix $\mathrm{pre}$, initialize the round machine with the encoded transition table,
the output summary and unary encoding of $q$, both auxiliary heads at zero, and
configuration data from $c$. There is a time $T$ with $0<T\le 50(L+1)$, where $L$ is the
length of the encoded transition table, at which the round machine reaches the
configuration built from applying $a$ to $c$, the new output summary and the unary
encoding of the successor state (empty if it halts), with output prefix $\mathrm{pre}$
and auxiliary heads at zero. At every strictly intermediate positive time its state is
neither the entry state nor the halt state, and at every time $u\le T$ its input head
and each original work head lie at either their starting or finishing positions.
\end{lemma}

\begin{definition}[Representation of a machine configuration]
\lean{Turing.StepBlock.Rep}
\label{Turing.StepBlock.Rep}
\leanok
\uses{Turing.Action, Turing.Cfg, Turing.FinTM.bufferTape, Turing.StepBlock.Ph, Turing.StepBlock.outputSummary, Turing.StepBlock.summaryCells, Turing.StepBlock.summaryTape, Turing.StepBlock.tableTape, Turing.tableBitsK, Turing.unaryFin}
Let $k,n\in\mathbb{N}$, let $x$ be a Boolean input word, let $\mathrm{tr}$ be a
transition function for a $k$-tape machine with Boolean work symbols and states in
$\text{Fin}(n+1)$, and let $c$ and $d$ be configurations of that machine and its
compiled machine, respectively. The representation condition requires that their
input-head positions agree; that the first $k$ work tapes and their head positions agree
cellwise; that the two additional tapes contain, respectively, the designated encoding
of the transition table of $\mathrm{tr}$ and the designated summary of $c$'s output and
state; and that both additional tape heads are at position zero.
\end{definition}

\begin{lemma}[Canonical Initial Round Configuration]
\lean{Turing.StepBlock.rep_entry}
\label{Turing.StepBlock.rep_entry}
\leanok
\uses{Turing.Action, Turing.Cfg, Turing.Cfg.mapState, Turing.StepBlock.Ph, Turing.StepBlock.Rep, Turing.StepBlock.atKey, Turing.StepBlock.outputSummary, Turing.StepBlock.roundData, Turing.StepBlock.summaryTape, Turing.StepBlock.tableTape, Turing.tableBitsK, Turing.unaryFin}
\difficulty{4}
Let $k,n\in\mathbb N$, let $x$ be a finite binary input, let $T$ be a transition
function for a machine with $k$ work tapes and states indexed by $\{0,\ldots,n\}$, let
$c$ be a configuration of that machine, let $q\in\{0,\ldots,n\}$, and let $d$ be a
configuration of the enlarged machine. Suppose the state of $c$ is $q$, the state of $d$
is the entry control state, and $d$ represents $c$ with respect to $T$. Then $d$ equals
the entry control state applied to the round configuration built from $c$, the
concatenated transition-table encoding, the output summary of $c$, the unary encoding of
$q$, zero positions for the table and state heads, and $d$'s output list.
\end{lemma}

\begin{theorem}[Round Machine Effect]
\lean{Turing.StepBlock.roundTM_effect}
\label{Turing.StepBlock.roundTM_effect}
\leanok
\uses{Turing.Action, Turing.Action.apply, Turing.Cfg, Turing.Cfg.inputSymbol, Turing.Cfg.mapState, Turing.Cfg.workTapeSymbols, Turing.MultiTapeTM.abs_workTapePos_runFrom_sub_le, Turing.MultiTapeTM.output_prefix, Turing.MultiTapeTM.runFrom, Turing.StepBlock.Ph, Turing.StepBlock.Rep, Turing.StepBlock.atKey, Turing.StepBlock.exit, Turing.StepBlock.outputSummary, Turing.StepBlock.outputSummary_append, Turing.StepBlock.rep_entry, Turing.StepBlock.roundData, Turing.StepBlock.roundTM, Turing.StepBlock.roundTM_canonical, Turing.StepBlock.summaryCells, Turing.StepBlock.summaryEmit, Turing.StepBlock.summaryTape, Turing.StepBlock.tableTape, Turing.moveInputPos, Turing.tableBitsK, Turing.workTapePos_apply_le}
\difficulty{6}
Let $k,n$ be natural numbers, let $x$ be a Boolean input word, let $M$ be a machine with
$k$ Boolean work tapes and $n+1$ states, and let $t$ be a natural number. Suppose a
configuration $c$ is in a live state and remains live through time $t$. Suppose also
that an augmented configuration $d$ is at the round machine's entry state and represents
$c$, including the encoded transition table and the auxiliary tapes. There is a time $T$
with $0<T\le 50(L+1)$, where $L$ is the encoded transition table length, such that the
round machine reaches a live entry state representing $M$'s one-step successor of $c$.
Its summary-tape symbol at that time is obtained by combining the initial summary-tape
symbol with the action's optional output: no output preserves the summary, and an
emitted bit makes the summary that bit if the old summary was absent, or $0$ otherwise.
At every strictly intermediate positive time the state is neither entry nor halt;
through time $T$ the output is unchanged, the input head is at its initial or final
position and moves by at most one from its initial position, and each original work head
is at its initial or final position. Every work-head displacement is at most $1$ on an
original tape and at most $50(L+1)$ on either auxiliary tape. Both auxiliary heads are
at zero initially and at time $T$.
\end{theorem}

\begin{lemma}[Visited cells over consecutive time blocks]
\lean{Turing.MultiTapeTM.visitedByTapeHead_add_subset}
\label{Turing.MultiTapeTM.visitedByTapeHead_add_subset}
\leanok
\uses{Turing.Cfg, Turing.MultiTapeTM, Turing.MultiTapeTM.runFrom, Turing.MultiTapeTM.runFrom_add, Turing.MultiTapeTM.visitedByTapeHead}
\difficulty{4}
Let $M$ be a deterministic multi-tape Turing machine with Boolean work-tape symbols and
state set $S$, let $x$ be a Boolean input word, and let $d$ be a configuration on $x$.
For natural numbers $T_1,T_2$ and a work-tape index $j$, the cells visited by tape $j$
from $d$ through step $T_1+T_2$ are contained in the union of the cells visited from $d$
through step $T_1$ and those visited from the configuration reached at step $T_1$
through step $T_2$.
\end{lemma}

\begin{lemma}[Visited cells after one initial step]
\lean{Turing.MultiTapeTM.visitedByTapeHead_step_subset}
\label{Turing.MultiTapeTM.visitedByTapeHead_step_subset}
\leanok
\uses{Turing.Cfg, Turing.MultiTapeTM, Turing.MultiTapeTM.runFrom_succ_eq_step, Turing.MultiTapeTM.runFrom_zero, Turing.MultiTapeTM.step, Turing.MultiTapeTM.visitedByTapeHead}
\difficulty{4}
Let $M$ be a deterministic multi-tape Turing machine with Boolean work-tape symbols and
state set $S$, let $x$ be a Boolean input word, and let $c$ be a configuration on $x$.
For every natural number $t$ and work-tape index $j$, the cell occupied by tape $j$'s
head in $c$, together with the cells visited by that tape from the one-step successor of
$c$ through step $t$, is contained in the cells visited from $c$ through step $t+1$.
\end{lemma}

\begin{theorem}[Simulation of one transition]
\lean{Turing.StepBlock.roundTM_visited}
\label{Turing.StepBlock.roundTM_visited}
\leanok
\uses{Turing.Action, Turing.Action.apply, Turing.Cfg, Turing.Cfg.inputSymbol, Turing.Cfg.workTapeSymbols, Turing.MultiTapeTM.runFrom, Turing.MultiTapeTM.visitedByTapeHead, Turing.StepBlock.Ph, Turing.StepBlock.Rep, Turing.StepBlock.exit, Turing.StepBlock.roundTM, Turing.StepBlock.roundTM_effect, Turing.tableBitsK}
\difficulty{4}
Let $k,n$ be nonnegative integers, let $x$ be a finite binary input, and let
$\mathrm{tr}$ be a transition rule for a machine with $k$ work tapes and states indexed
by $\{0,\ldots,n\}$. Suppose a configuration $c$ is in state $q$, and let $d$ be a
configuration of the machine that simulates one transition, with its state at the entry
phase and its work tapes and head positions related to those of $c$ as prescribed by the
representation relation. Write $a$ for the action selected by $\mathrm{tr}$ from state
$q$, the symbol under the input head of $c$, and the symbols under its work-tape heads.
Then there is an integer $T$ with $0<T\le 50(L+1)$, where $L$ is the length of the
binary encoding of the transition table, such that after $T$ steps the simulating
machine is at the exit phase determined by whether $a$ has a successor state, and its
configuration represents the result of applying $a$ to $c$. At every intermediate time
$u$ with $0<u<T$, its state is neither the entry nor the halt anchor; its output is
unchanged through time $T$. Each original work-tape head visits only its position in $c$
or its position after applying $a$, and each of the two additional work-tape heads
visits only positions in the interval $[-50(L+1),50(L+1)]$.
\end{theorem}

\begin{theorem}[Iteration of the Round Machine]
\lean{Turing.StepBlock.roundTM_iterate}
\label{Turing.StepBlock.roundTM_iterate}
\leanok
\uses{Turing.Action.apply, Turing.Cfg, Turing.Cfg.inputSymbol, Turing.Cfg.workTapeSymbols, Turing.MultiTapeTM, Turing.MultiTapeTM.runFrom, Turing.MultiTapeTM.runFrom_add, Turing.MultiTapeTM.runFrom_succ_eq_step, Turing.MultiTapeTM.runFrom_zero, Turing.MultiTapeTM.step, Turing.MultiTapeTM.visitedByTapeHead, Turing.MultiTapeTM.visitedByTapeHead_add_subset, Turing.MultiTapeTM.visitedByTapeHead_step_subset, Turing.StepBlock.Ph, Turing.StepBlock.Rep, Turing.StepBlock.roundTM, Turing.StepBlock.roundTM_visited, Turing.StepBlock.summaryTape, Turing.StepBlock.tableTape, Turing.tableBitsK}
\difficulty{6}
Let $k,n,t$ be natural numbers, let $x$ be a Boolean input word, and let $M$ be a
machine with $k$ Boolean work tapes and $n+1$ states. Suppose $c$ remains in a live
state through time $t$. For every augmented configuration $d$ at the round machine's
entry state that represents $c$ and $M$'s transition rule, there is a time $T$ with
$t\le T\le t\,50(L+1)$, where $L$ is the encoded transition table length, such that the
round machine is again at its entry state and represents $M$'s configuration at time
$t$. Through time $T$, its output remains unchanged and it never reaches the halt state.
Every original work-tape cell visited by the round machine's corresponding head through
time $T$ is visited by the corresponding head of $M$ through time $t$; every cell
visited by either auxiliary head lies in the integer interval $[-50(L+1),50(L+1)]$.
\end{theorem}

\begin{theorem}[Unique First Exit Time]
\lean{Turing.StepBlock.roundTM_firstExit_unique}
\label{Turing.StepBlock.roundTM_firstExit_unique}
\leanok
\uses{Turing.Action, Turing.Cfg, Turing.MultiTapeTM.runFrom, Turing.StepBlock.Ph, Turing.StepBlock.Rep, Turing.StepBlock.exit, Turing.StepBlock.roundTM, Turing.StepBlock.roundTM_effect}
\difficulty{5}
Let $k,n$ be natural numbers, let $x$ be a Boolean input word, let $\mathrm{tr}$ be a
transition rule for $n+1$ states, optional input bits, and $k$ work-tape reads, and let
$c$ be a configuration with state $q$. Let $d$ be an augmented configuration at the
round machine's entry state that represents $c$ and the transition rule. There is a
unique positive time $T$ at which the round machine is in either its entry state or its
halt state; at every strictly intermediate positive time it is in neither state.
\end{theorem}
